% concatenated from main.tex
\documentclass[10pt]{amsart} 
\usepackage{amssymb}
\font\grd cmr10 scaled\magstep2 \font\grgd cmbx10 scaled\magstep2
\def\demo{\noindent{\bf Proof .}}
\newtheorem{thm1}{Theorem}[section]
\newtheorem{lem1}[thm1]{Lemma}
\newtheorem{rem1}[thm1]{Remark}
\newtheorem{def1}[thm1]{Definition}
\newtheorem{def2}[thm1]{Definitions}
\newtheorem{cor1}[thm1]{Corollary}
\newtheorem{defn1}[thm1]{Definition}
\newtheorem{prop1}[thm1]{Proposition}
\newtheorem{ex1}[thm1]{Example}
\newtheorem{not1}[thm1]{Notation}

%\DeclareMathOperator{\deg}{deg}
\DeclareMathOperator{\supp}{supp} \DeclareMathOperator{\lcm}{lcm}
\DeclareMathOperator{\ind}{ind}

% Template article for preprint document class `elsart'
\begin{document}

\title[]
{On complete intersection projective closures of monomial curves}
\author[]{Anargyros Katsabekis}
\address {Department of Mathematics, University of Ioannina, 45110 Ioannina, Greece} \email{katsampekis@uoi.gr}
\keywords{Complete intersection; Arithmetically Cohen--Macaulay; Projective closure; Monomial curve}
\subjclass{13F65, 13H10, 13P10, 14M10, 14M25}




\begin{abstract} In this paper, we study the complete intersection property of projective closures of affine monomial curves. For monomial curves in the four-dimensional affine space, we give a complete and explicit characterization of when the defining ideals of their projective closures are complete intersections. We also investigate when the projective closures are arithmetically Cohen--Macaulay.
\end{abstract}
\maketitle





\section{Introduction}
%-------------------------------------------------

Toric varieties are fundamental objects in algebraic geometry and commutative algebra. Among their key properties, the complete intersection property is particularly important, as it implies that the associated coordinate ring is Cohen--Macaulay. A toric variety is called a {\em complete intersection} if its associated toric ideal is a {\em complete intersection}, that is, if it can be generated by a number of polynomials equal to its height.

Determining when the projective closure of a toric variety is a complete intersection requires a detailed analysis, since homogenization of the associated toric ideal may increase the minimal number of generators and alter the defining equations. In this paper, we investigate when the projective closure of an affine monomial curve is a complete intersection. While the complete intersection property has been thoroughly studied for affine monomial curves (see, e.g., \cite{Bresinsky75}, \cite{Her}), much less is known about
their projective closures, precisely because homogenization can significantly
change the structure of the defining toric ideal. Note that the projective closure can be a complete intersection only if the affine monomial curve itself is. Indeed, dehomogenizing a generating set of the defining ideal of the projective closure yields a generating set of the affine toric ideal.


The paper \cite{BG2} studies the complete intersection property for projective closures of certain families of monomial curves in the four-dimensional affine space, while \cite{BG} treats the analogous question for simplicial projective toric varieties. Despite these contributions, no complete characterization has been known for projective closures of complete intersection affine monomial curves in the four-dimensional affine space. Our main result gives necessary and sufficient
conditions for the projective closure $\overline{C({\bf a})}$ to be a
complete intersection, expressed explicitly in terms of the minimal binomial generators
of the affine toric ideal $I({\bf a})$. Our classification covers all three structural cases of Theorem~\ref{CasesMinimal}, whereas \cite{BG2} addresses only certain
families within these cases. Our approach combines a detailed analysis of
these generators with Gr\"obner basis techniques.

We also investigate the arithmetically Cohen--Macaulay property of the projective closure, providing a complete characterization in Case~1 and partial results in Case~2; the full characterization in Cases~2 and~3 remains
an open problem. In addition, Proposition~\ref{ACM7} constructs infinite families of monomial 
curves in the four-dimensional affine space whose projective closures are arithmetically 
Cohen--Macaulay but whose homogeneous toric ideals are not complete 
intersections and have arbitrarily many minimal generators. In particular, 
the arithmetically Cohen--Macaulay property does not imply the complete 
intersection property.


Let ${\bf a}=(a_{1},\ldots,a_{n})$ be a sequence of distinct positive integers with ${\rm gcd}(a_{1},\ldots,a_{n})=1$, and let $C({\bf a})$ be the affine monomial curve parametrized by $x_{i}=t^{a_i}$ for $1 \leq i \leq n$. Consider the polynomial ring $K[x_{1},\ldots,x_{n}]$ over any field $K$. The {\em toric ideal} of $C({\bf a})$, denoted by $I({\bf a})$, is the kernel of the $K$-algebra homomorphism $\phi \colon K[x_{1},\ldots,x_{n}] \to K[t]$ given by $\phi(x_{i})=t^{a_i}$ for all $1 \leq i \leq n$. By \cite[Lemma 4.1]{Sturmfels95}, $I({\bf a})$ is generated by binomials of the form $x_{1}^{u_1} \cdots x_{n}^{u_n}-x_{1}^{v_1} \cdots x_{n}^{v_n}$ such that $u_{1}a_{1}+\cdots+u_{n}a_{n}=v_{1}a_{1}+\cdots+v_{n}a_{n}$. Moreover, $I({\bf a})$ contains no binomial of the form $x_{1}^{u_1} \cdots x_{n}^{u_n}-1$ with $(u_1,\ldots,u_n) \neq (0,\ldots,0)$. By \cite[Lemma 4.2]{Sturmfels95}, the height of $I({\bf a})$ is $n-1$, and hence
$I({\bf a})$ is a complete intersection if and only if it is minimally generated by $n-1$ binomials.

Set $d={\rm max}\{a_1,\ldots,a_n\}$. The {\em homogenization} $I^{h}({\bf a})$ of $I({\bf a})$ with respect to the variable $x_{0}$ is the kernel of the $K$-algebra homomorphism $\psi \colon K[x_{0},x_1,\ldots,x_n] \to K[s,t]$, defined by $\psi(x_{0})=s^{d}$ and $\psi(x_{i})=s^{d-a_{i}}t^{a_i}$ for all $1 \leq i \leq n$. Then $I^{h}({\bf a})$ is a complete intersection if and only if it is minimally generated by $n-1$ binomials. 

The {\em projective closure} $\overline{C({\bf a})} \subset \mathbb{P}^{n}$ of $C({\bf a})$ is the projective monomial curve parametrized by $[s:t] \mapsto [s^{d}:s^{d-a_{1}}t^{a_1}: \cdots: s^{d-a_{n}}t^{a_n}]$ and its defining homogeneous ideal is $I^{h}({\bf a})$. The curve $\overline{C({\bf a})}$ is called {\em arithmetically Cohen--Macaulay} if $I^{h}({\bf a})$ is Cohen--Macaulay. This property has been extensively studied; see, for example, \cite{HerSta}, \cite{Katsa1}, \cite{PatilRoberts2003}, and \cite{SSS}.

In Section 2, we observe that when $n=4$ and $I({\bf a})$ is a complete intersection toric ideal, its minimal generating set falls into one of three distinct cases (see Theorem \ref{CasesMinimal}). For the first case, we establish necessary and sufficient conditions for the projective closure to be arithmetically Cohen–Macaulay (see Theorems \ref{ACM1}, \ref{ACM2}, \ref{ACM3}, and \ref{ACM4}) and for $I^{h}({\bf a})$ to be a complete intersection (see Theorems \ref{BAS1} and \ref{CI23}, Corollaries \ref{CorCase11CI} and \ref{CorCase1CI}).

In Section 3, we consider the second case and determine necessary and sufficient conditions under which $I^{h}({\bf a})$ is a complete intersection (see Theorems \ref{Veryfirst}, \ref{Basic4}, \ref{Comp2}, and Corollaries \ref{Cor12}, \ref{Cor13}). We also establish necessary and sufficient conditions for $\overline{C({\bf a})}$ to be arithmetically Cohen--Macaulay in certain subcases (see Theorems \ref{ACM5}, \ref{ACM6}), and construct families of complete intersection affine monomial curves $C({\bf a})$ whose projective closures are arithmetically Cohen--Macaulay, even though the ideals $I^{h}({\bf a})$ have a large number of minimal generators (see Proposition \ref{ACM7}).

In Section 4, we study the third case and derive explicit criteria for $I^{h}({\bf a})$ to be a complete intersection (see Theorems \ref{Basic5}, \ref{Basic23}, \ref{basicvfff}, \ref{Basic24}, \ref{Basic25}, \ref{Basic26}).



\section{Projective closures in case 1}

In this section, we study in detail the first case arising from the structure of the minimal generators. For any $f \in K[x_1,\ldots,x_n]$, we denote by $f^{h}$ its homogenization with respect to $x_0$. The following remark will be used repeatedly.

\begin{rem1}\label{VeryBasicRemark} {\rm \begin{enumerate} \item Let $I\subset K[x_1,\ldots,x_n]$ be an ideal generated by $f_1,\ldots,f_r$. If the initial monomials of the generators
$f_1,\ldots,f_r$
are pairwise relatively prime, then $\{f_1,\ldots,f_r\}$ is a Gr\"obner basis for $I$ with respect to $<$ by Buchberger's criterion \cite[Proposition~4, Chapter~2]{CLO}. \item If $\{f_1,\ldots,f_r\}$ is a Gr\"obner basis for an ideal $I$ with respect to a graded monomial order $<$, then $\{f_1^h,\ldots,f_r^h\}$ is a Gr\"obner basis for $I^h$ with respect to the graded monomial order $<_0$ extending $<$ and satisfying $x_0<_0x_i$ for all $i\ge1$ (see \cite[Chapter~8]{CLO}). \end{enumerate} } \end{rem1}

Let ${\bf a}=(a_1,\ldots,a_4)$ be a sequence of distinct relatively prime positive integers. 
\begin{not1} {\rm Throughout the paper, let $c_i$ denote the smallest positive integer such that $c_{i}a_{i} \in \sum_{j \neq i} \mathbb{N}a_{j}$, for each $i=1,\ldots,4$, where $\mathbb{N}$ is the set of nonnegative integers.}
\end{not1}

\begin{thm1} \label{CasesMinimal} Let $I({\bf a})$ be a complete intersection toric ideal, where ${\bf a}=(a_1,\ldots,a_4)$ is a sequence of distinct positive integers with ${\rm gcd}(a_1,\ldots,a_4)=1$. Then after permuting the variables, if necessary, there exists a minimal system of binomial generators $S$ of $I({\bf a})$ of the following form \begin{enumerate} \item[{\rm CASE 1:}] If $c_{1}a_{1}=c_{2}a_{2}$ and $c_{3}a_{3}=c_{4}a_{4}$, then $S=\{x_{2}^{c_2}-x_{1}^{c_1}, x_{3}^{c_3}-x_{4}^{c_4}, x_{1}^{d_1}x_{2}^{d_2}-x_{3}^{d_3}x_{4}^{d_4}\}$. \item[{\rm CASE 2:}] If $c_{1}a_{1}=c_{2}a_{2}=c_{3}a_{3} \neq c_{4}a_{4}$, then $S=\{x_{2}^{c_2}-x_{1}^{c_1}, x_{3}^{c_3}-x_{1}^{c_1}, x_{4}^{c_4}-x_{1}^{d_1}x_{2}^{d_2}x_{3}^{d_3}\}$, where at least two of $d_1$, $d_2$, and $d_3$ are nonzero. \item[{\rm CASE 3:}] If $c_{1}a_{1}=c_{2}a_{2}$ and $c_{i}a_{i} \neq c_{j}a_{j}$ for all $\{i,j\} \neq \{1,2\}$, then $S=\{x_{2}^{c_2}-x_{1}^{c_1}, x_{3}^{c_3}-x_{1}^{b_1}x_{2}^{b_2}, x_{4}^{c_4}-x_{1}^{d_1}x_{2}^{d_2}x_{3}^{d_3}\}$, where both $b_1$ and $b_2$ are nonzero, and at least two of $d_1$, $d_2$, and $d_3$ are nonzero.
\end{enumerate}
\end{thm1}

{\em \noindent Proof.} We begin by noting that if $I({\bf a})$ is a complete intersection, then one of the following holds: either $c_1a_1=c_2a_2$ and $c_3a_3=c_4a_4$; or $c_{1}a_{1}=c_{2}a_{2}=c_{3}a_{3} \neq c_{4}a_{4}$; or $c_{1}a_{1}=c_{2}a_{2}$ and $c_{i}a_{i} \neq c_{j}a_{j}$ for all $\{i,j\} \neq \{1,2\}$ (see \cite{KO}). By \cite[Proposition 3.2]{Shi}, a minimal generating set of $I({\bf a})$ is of the form $\{x_{2}^{c_2}-x_{1}^{c_1}, x_{3}^{c_3}-x_{4}^{c_4}, x_{1}^{d_1}x_{2}^{d_2}-x_{3}^{d_3}x_{4}^{d_4}\}$ in Case 1, and $\{x_{2}^{c_2}-x_{1}^{c_1}, x_{3}^{c_3}-x_{1}^{b_1}x_{2}^{b_2}, x_{4}^{c_4}-x_{1}^{d_1}x_{2}^{d_2}x_{3}^{d_3}\}$ in Cases 2 and 3. In Case 2, since $c_{1}a_1=c_{2}a_2=c_{3}a_{3}$, we can set $b_{1}=c_1$ and $b_2=0$, so that the second generator is $x_{3}^{c_3}-x_{1}^{c_1}$. We may also assume that at least two of $d_1$, $d_2$, $d_3$ are nonzero; otherwise, if, for instance, $d_{1}=d_{2}=0$, then $d_{3}>c_3$, and using $x_{3}^{c_3}-x_{1}^{c_1}$, we can replace $x_{4}^{c_4}-x_{3}^{d_3}$ with $x_{4}^{c_4}-x_{1}^{c_1}x_{3}^{d_3-c_3} \in I({\bf a})$. In Case 3, we may assume that both $b_1$ and $b_2$ are nonzero; indeed, if, for instance, $b_2=0$, then $b_1>c_1$, and using $x_2^{c_2}-x_1^{c_1}$ we can replace $x_3^{c_3} - x_1^{b_1}x_2^{b_2}$ with $x_3^{c_3}-x_1^{b_1-c_1}x_2^{b_2+c_2} \in I({\bf a})$. Similarly, we may assume that at least two of $d_1$, $d_2$, $d_3$ are nonzero. \hfill $\square$ \par\vspace{0.5\baselineskip}

 For the rest of this section, let $I({\bf a})$ be a complete intersection toric ideal minimally generated by $x_{2}^{c_2}-x_{1}^{c_1}$, $x_{3}^{c_3}-x_{4}^{c_4}$, and $x_{1}^{d_1}x_{2}^{d_2}-x_{3}^{d_3}x_{4}^{d_4}$. {\em Without loss of generality, we fix the assumptions $a_{4}={\rm max}\{a_1,\ldots,a_4\}$, $c_2>c_1$, $d_2 \leq c_2$, and $d_3 \leq c_3$.} Indeed, if $d_2>c_2$, we may write $d_{2}=kc_2+l$ for some integers $k \geq 1$ and $0 \leq l<c_2$, and replace $x_{1}^{d_1}x_{2}^{d_2}-x_{3}^{d_3}x_{4}^{d_4}$ with $x_{1}^{kc_1+d_1}x_{2}^{l}-x_{3}^{d_3}x_{4}^{d_4} \in I({\bf a})$. Similarly, we may assume $d_3 \leq c_3$. Finally, since $a_4>a_3$, we have $c_4<c_3$.

Throughout this section, let $<$ be the degree reverse lexicographic order on $K[x_1, \ldots, x_4]$ with $x_1 > x_2 > x_3 > x_4$. We begin by considering a first subcase in which the ideal $I^{h}({\bf a})$ is a complete intersection. \vspace{-0.25\baselineskip}

\begin{thm1} \label{BAS1} If $d_2=c_2$ and $d_3=c_3$, then $I^{h}({\bf a})$ is a complete intersection.
\end{thm1}
\vspace{-0.25\baselineskip}
{\em \noindent Proof.}The binomials
\[
x_2^{c_2}-x_1^{c_1},\qquad
x_3^{c_3}-x_4^{c_4},\qquad
x_1^{c_1+d_1}-x_4^{c_4+d_4}
\]
belong to $I({\bf a})$.
Their initial monomials are
$x_2^{c_2}$,
$x_3^{c_3}$,
and
$x_1^{c_1+d_1}$,
which are pairwise relatively prime.
Hence, by Remark~\ref{VeryBasicRemark}(1), these binomials form a
Gr\"obner basis for $I({\bf a})$ with respect to $<$.
By Remark~\ref{VeryBasicRemark}(2), their homogenizations form a
Gr\"obner basis for $I^h({\bf a})$ with respect to $<_0$.
Hence $I^h({\bf a})$ is a complete intersection.
\hfill $\square$ \par\vspace{0.5\baselineskip}

The following example illustrates the subcase considered in Theorem~\ref{BAS1}.

\begin{ex1} {\rm Let ${\bf a}=(132, 60, 165, 220)$. Then $I({\bf a})$ is generated by $x_{2}^{11}-x_{1}^{5}$, $x_3^{4}-x_{4}^{3}$, and $x_{2}^{11}-x_{3}^{4}$. Here, $d_{2}=c_2=11$ and $d_{3}=c_3=4$. By Theorem \ref{BAS1}, the ideal $I^{h}({\bf a})$ is a complete intersection.}
\end{ex1}

Next, we consider another subcase in which $I^{h}({\bf a})$ is a complete intersection.

\begin{thm1} \label{CI23} Suppose that $d_2=0$ and $d_{3}=c_{3}$. Then $I^{h}({\bf a})$ is a complete intersection.
\end{thm1}
{\em \noindent Proof.} Let
$f_1=x_2^{c_2}-x_1^{c_1}$,
$f_2=x_3^{c_3}-x_4^{c_4}$,
and
$g=x_1^{d_1}-x_4^{c_4+d_4}\in I({\bf a})$,
since $d_2=0$ and $d_3=c_3$.
The initial monomials of $f_1$, $f_2$, and $g$ are
$x_2^{c_2}$,
$x_3^{c_3}$,
and
$x_1^{d_1}$,
respectively, and are pairwise relatively prime.
Hence, by Remark~\ref{VeryBasicRemark}(1),
$\{f_1,f_2,g\}$ is a Gr\"obner basis for $I({\bf a})$ with respect to $<$.
By Remark~\ref{VeryBasicRemark}(2),
$\{f_1^{h},f_2^{h},g^{h}\}$ is a Gr\"obner basis for $I^{h}({\bf a})$
with respect to $<_0$.
Thus $I^{h}({\bf a})$ is a complete intersection.
\hfill $\square$ \par\vspace{0.5\baselineskip}

We now obtain several sufficient and necessary conditions under which the arithmetically Cohen--Macaulay property holds.

\begin{thm1} \label{ACM1} Suppose that $d_2=c_2$ and $d_3<c_3$. Then $\overline{C({\bf a})}$ is arithmetically Cohen--Macaulay if and only if $c_{1}+d_{1} \geq d_{3}+d_{4}$.
\end{thm1}

{\em \noindent Proof.} ($\Leftarrow$) Let $f_{1}=x_2^{c_2}-x_1^{c_1}$, $f_{2}=x_3^{c_3}-x_4^{c_4}$, and $f_{3}=x_1^{c_1+d_{1}}-x_3^{d_3}x_4^{d_4}\in I({\bf a})$. Then ${\rm in}_{<}(f_{1})=x_2^{c_2}$, ${\rm in}_{<}(f_{2})=x_3^{c_3}$, and ${\rm in}_{<}(f_{3})=x_1^{c_1+d_1}$ are pairwise relatively prime. Hence, by Remark~\ref{VeryBasicRemark}(1), $\{f_1,f_2,f_3\}$ is a Gr\"obner basis for $I({\bf a})$ with respect to $<$. Since none of these initial monomials is divisible by $x_4$, it follows from \cite[Theorem~2.2]{HerSta} that $\overline{C({\bf a})}$ is arithmetically Cohen--Macaulay.\\
($\Rightarrow$) Assume that $c_{1}+d_{1}<d_{3}+d_{4}$. Since $\overline{C({\bf a})}$ is arithmetically Cohen--Macaulay and ${\rm in}_{<}(f_{3})=x_3^{d_{3}}x_4^{d_4}$, by \cite[Theorem 2.2]{HerSta} there exists a binomial $B \in I({\bf a})$ such that ${\rm in}_{<}(B)$ divides $x_3^{d_3}$; in particular ${\rm in}_{<}(B)=x_{3}^{e}$ for some $e \leq d_3$. Since $d_3<c_3$, this gives $e<c_3$, contradicting the minimality of $c_3$. \hfill $\square$ \par\vspace{0.5\baselineskip}

As an application of Theorem \ref{ACM1}, we obtain the following corollary. 

\begin{cor1} \label{CorCase11CI} Suppose that $d_2=c_2$ and $d_{3}<c_{3}$. Then the ideal $I^{h}({\bf a})$ is a complete intersection if and only if $c_{1}+d_{1} \geq d_{3}+d_{4}$.
\end{cor1}

{\em \noindent Proof.} ($\Leftarrow$) By the proof of Theorem~\ref{ACM1}, $\{f_i \mid 1\le i\le 3\}$ is a Gr\"obner basis for $I({\bf a})$ with respect to $<$. Hence, by Remark~\ref{VeryBasicRemark}(2), $\{f_i^{h} \mid 1\le i\le 3\}$ is a Gr\"obner basis for $I^{h}({\bf a})$ with respect to $<_0$. Thus $I^{h}({\bf a})$ is a complete intersection.\\ ($\Rightarrow$) If $I^{h}({\bf a})$ is a complete intersection, then $\overline{C({\bf a})}$ is arithmetically Cohen--Macaulay. By Theorem \ref{ACM1}, it follows that $c_1+d_{1} \geq d_{3}+d_{4}$. \hfill $\square$ \par\vspace{0.5\baselineskip}


\begin{thm1} \label{ACM2} Suppose that $0<d_2<c_2$ and $d_{3}=c_{3}$. Then $\overline{C({\bf a})}$ is arithmetically Cohen--Macaulay if and only if $c_{1}+d_{1} \geq c_{2}-d_{2}+c_{4}+d_{4}$.
\end{thm1}

{\em \noindent Proof.} ($\Leftarrow$) Consider the set $T=\{f_1=x_2^{c_2}-x_1^{c_1},\; f_2=x_3^{c_3}-x_4^{c_4},\; f_3=x_1^{d_1}x_2^{d_2}-x_4^{c_4+d_4},\; f_4=x_1^{c_1+d_1}-x_2^{c_2-d_2}x_4^{c_4+d_4}\} \subset I({\bf a})$. Since $c_2>c_1$ and $c_1+d_1 \geq c_2-d_2+c_4+d_4$, the initial monomials with respect to $<$ are ${\rm in}_{<}(f_{1})=x_2^{c_2}$, ${\rm in}_{<}(f_{2})=x_3^{c_3}$, ${\rm in}_{<}(f_{3})=x_1^{d_1}x_2^{d_2}$, and ${\rm in}_{<}(f_4)=x_1^{c_1+d_{1}}$. Since ${\rm in}_{<}(f_1)$ and ${\rm in}_{<}(f_2)$ are relatively prime, it follows that $S(f_{1},f_{2}) \stackrel{T}{\longrightarrow} 0$. Similarly, $S(f_{1},f_4) \stackrel{T}{\longrightarrow} 0$, $S(f_{2},f_{3}) \stackrel{T}{\longrightarrow} 0$, and $S(f_{2},f_4) \stackrel{T}{\longrightarrow} 0$. We have that $S(f_{1},f_{3})=-f_4 \stackrel{f_4}{\longrightarrow} 0$ and $S(f_{3},f_4)=x_4^{c_4+d_4}(x_{2}^{c_2}-x_{1}^{c_1}) \stackrel{f_1}{\longrightarrow} 0$. Thus $T$ is a Gr\"obner basis for $I({\bf a})$ with respect to $<$. Since none of the initial monomials is divisible by $x_4$, it follows from \cite[Theorem~2.2]{HerSta} that $\overline{C(\mathbf a)}$ is arithmetically Cohen--Macaulay.\\
($\Rightarrow$) Assume that $c_{1}+d_{1}<c_{2}-d_{2}+c_{4}+d_{4}$. Since $\overline{C({\bf a})}$ is arithmetically Cohen--Macaulay and ${\rm in}_{<}(f_4)=x_2^{c_2-d_2}x_4^{c_4+d_4}$, by \cite[Theorem~2.2]{HerSta} there exists a binomial $B \in I({\bf a})$ such that ${\rm in}_{<}(B)$ divides $x_2^{c_{2}-d_2}$, contradicting the minimality of $c_2$. \hfill $\square$


\begin{thm1} \label{ACM3} Suppose that $0<d_2<c_2$ and $d_{3}<c_{3}$. Then $\overline{C({\bf a})}$ is arithmetically Cohen--Macaulay if and only if $c_{1}+d_{1} \geq c_{2}-d_{2}+d_{3}+d_{4}$.
\end{thm1}

{\em \noindent Proof.} ($\Leftarrow$) The proof proceeds via Gr\"obner basis techniques, as in the proof of Theorem \ref{ACM2}. Since $c_2 > c_1$ and $c_1 + d_1 \geq c_2 - d_2 + d_3 + d_4$, it follows that $d_1 + d_2 > d_3 + d_4$. Let $f_1=x_2^{c_2}-x_1^{c_1}, f_2=x_3^{c_3}-x_4^{c_4}, f_3=x_1^{d_1}x_2^{d_2}-x_3^{d_3}x_4^{d_4}, f_4=x_1^{c_1+d_{1}}-x_2^{c_{2}-d_2}x_3^{d_3}x_4^{d_4} \in I({\bf a})$. With respect to $<$, the initial monomials of the above binomials are $x_2^{c_2}$, $x_3^{c_3}$, $x_1^{d_1}x_2^{d_2}$, and $x_1^{c_1+d_{1}}$, respectively. As in the proof of Theorem~\ref{ACM2}, all $S$-polynomials reduce to zero, so $\{f_i \mid 1 \leq i \leq 4\}$ is a Gr\"obner basis of $I(\mathbf a)$ with respect to $<$. Since $x_4$ does not divide any initial monomial, $\overline{C({\bf a})}$ is arithmetically Cohen--Macaulay by \cite[Theorem 2.2]{HerSta}.\\
($\Rightarrow$) Assume that $c_{1}+d_{1}<c_{2}-d_{2}+d_{3}+d_{4}$. Since $\overline{C({\bf a})}$ is arithmetically Cohen--Macaulay and ${\rm in}_{<}(f_4)=x_2^{c_{2}-d_2}x_3^{d_3}x_4^{d_4}$, by \cite[Theorem 2.2]{HerSta} there exists a binomial $B \in I({\bf a})$ such that ${\rm in}_{<}(B)$ divides $x_2^{c_{2}-d_2}x_3^{d_3}$; in particular ${\rm in}_{<}(B)=x_2^{a}x_3^{b}$ for some $a \leq c_{2}-d_{2}$ and $b \leq d_{3}$. Note that $a \neq 0$. If $a = 0$, then ${\rm in}_{<}(B)= x_3^b$ with $b < c_3$, contradicting the minimality of $c_3$. Similarly, $b \neq 0$. Moreover, $d_1 \neq 0$ since $d_2 < c_2$, and $d_4 \neq 0$ since $d_3<c_3$. Since $B \in I({\bf a})$ and $I({\bf a})$ is generated by the binomials $f_1$, $f_2$, and $f_3$, the monomial $x_2^{a}x_3^{b}$ must be divisible by one of the monomials $x_2^{c_2}$, $x_3^{c_3}$, $x_1^{d_1}x_2^{d_2}$, or $x_3^{d_3}x_4^{d_4}$, a contradiction. \hfill $\square$

\begin{rem1} {\rm In each of the cases considered in Theorems \ref{ACM2} and \ref{ACM3}, $\overline{C({\bf a})}$ is arithmetically Cohen--Macaulay. Moreover, in each case $I^{h}({\bf a})$ admits a Gr\"obner basis with more than three elements. Thus $I^{h}({\bf a})$ is not a complete intersection.}
\end{rem1}


\begin{thm1} \label{ACM4} Suppose that $d_2=0$ and $d_{3}<c_{3}$. Then $\overline{C({\bf a})}$ is arithmetically Cohen--Macaulay if and only if $d_{1} \geq d_{3}+d_{4}$.
\end{thm1}
{\em \noindent Proof.} ($\Leftarrow$) Let $f_{1}=x_2^{c_2}-x_1^{c_1}, f_{2}=x_3^{c_3}-x_4^{c_4}$, and $f_{3}=x_1^{d_1}-x_3^{d_3}x_4^{d_4}$. Then their initial monomials $x_2^{c_2}$, $x_3^{c_3}$, and $x_1^{d_1}$ are pairwise relatively prime, hence $\{f_i \mid 1 \leq i \leq 3 \}$ is a Gr\"obner basis of $I({\bf a})$ with respect to $<$ by Remark~\ref{VeryBasicRemark}(1). Since none of the initial monomials is divisible by $x_4$, $\overline{C({\bf a})}$ is arithmetically Cohen--Macaulay by \cite[Theorem 2.2]{HerSta}.\\
($\Rightarrow$) Assume that $d_{1}<d_{3}+d_{4}$. Since $\overline{C({\bf a})}$ is arithmetically Cohen--Macaulay and ${\rm in}_{<}(f_{3})=x_3^{d_3}x_4^{d_4}$, \cite[Theorem 2.2]{HerSta} yields a binomial $B \in I({\bf a})$ such that ${\rm in}_{<}(B)$ divides $x_3^{d_{3}}$, contradicting the minimality of $c_{3}$. \hfill $\square$ \par\vspace{0.5\baselineskip}

As an application of Theorem \ref{ACM4}, we obtain the following corollary. 

\begin{cor1} \label{CorCase1CI} Suppose that $d_2=0$ and $d_{3}<c_{3}$. Then the ideal $I^{h}({\bf a})$ is a complete intersection if and only if $d_{1} \geq d_{3}+d_{4}$.
\end{cor1}

{\em \noindent Proof.} ($\Leftarrow$) By the proof of Theorem~\ref{ACM4}, $\{f_i \mid 1\le i\le 3\}$ is a Gr\"obner basis for $I({\bf a})$ with respect to $<$. Hence, by Remark~\ref{VeryBasicRemark}(2), $\{f_i^{h} \mid 1\le i\le 3\}$ is a Gr\"obner basis for $I^{h}({\bf a})$ with respect to $<_0$. Thus $I^{h}({\bf a})$ is a complete intersection.\\ ($\Rightarrow$) If $I^{h}({\bf a})$ is a complete intersection, then $\overline{C({\bf a})}$ is arithmetically Cohen--Macaulay. By Theorem \ref{ACM4}, it follows that $d_{1} \geq d_{3}+d_{4}$. \hfill $\square$ \par\vspace{0.5\baselineskip}

We conclude this section with an example illustrating the complete intersection criterion given in Corollary~\ref{CorCase1CI}.

\begin{ex1} {\rm Let ${\bf a}=(14, 7, 8, 20)$ and ${\bf v}=(7,7,4,10)$. By \cite[Theorem 2.6]{Katsa}, for any $t \in \mathbb{N}$ such that ${\rm gcd}(14+7t,7+7t,8+4t,20+10t)=1$, the ideal $I({\bf a}+t{\bf v})$ is generated by $x_{2}^{2+t}-x_{1}^{1+t}$, $x_{3}^{5}-x_{4}^{2}$, and $x_{1}^{2}-x_{3}x_{4}$. Here, $d_1=2$, $d_2=0$, $d_3=d_4=1$, and $c_3=5$. By Corollary \ref{CorCase1CI}, the ideal $I^{h}({\bf a}+t{\bf v})$ is a complete intersection whenever ${\rm gcd}(14+7t,7+7t,8+4t,20+10t)=1$.}
\end{ex1}




\section{Projective closures in case~2}

This section focuses on Case 2 of Theorem \ref{CasesMinimal}, where $I({\bf a})$ has a minimal generating set $\{g_{1}=x_{2}^{c_2}-x_{1}^{c_1}, g_{2}=x_{3}^{c_3}-x_{1}^{c_1}, g_{3}=x_{4}^{c_4}-x_{1}^{d_1}x_{2}^{d_2}x_{3}^{d_3}\}$, and at least two of $d_1$, $d_2$, and $d_3$ are nonzero.

We begin by studying the case where $a_{i} = {\rm max}\{a_1,\ldots,a_4\}$ for some $1 \leq i \leq 3$. {\em Without loss of generality, assume $a_1={\rm max}\{a_1,\ldots,a_4\}$, and fix $d_{2} \leq c_2$ and $d_{3} \leq c_3$.}

Throughout this section, unless stated otherwise, let $<$ denote the degree
reverse lexicographic order on $K[x_1,\ldots,x_4]$ with
$x_4>x_3>x_2>x_1$.

We first consider the case $d_2=c_{2}$ and $d_{3}=c_3$.

\begin{thm1} \label{Veryfirst} If $d_2=c_{2}$ and $d_{3}=c_3$, then $I^{h}({\bf a})$ is a complete intersection.
\end{thm1}

{\em \noindent Proof.} Using $g_1$ and $g_2$, we obtain $f=x_{4}^{c_4}-x_{1}^{d_1+2c_1} \in I({\bf a})$. Then ${\rm in}_{<}(g_1)=x_{2}^{c_2}$, ${\rm in}_{<}(g_2)=x_{3}^{c_3}$, and ${\rm in}_{<}(f)=x_{4}^{c_4}$ are pairwise relatively prime. Hence, by Remark~\ref{VeryBasicRemark}(1), $\{g_1, g_2, f\}$ is a Gr\"obner basis for $I({\bf a})$ with respect to $<$. By Remark \ref{VeryBasicRemark}(2), $\{g_{1}^{h}, g_{2}^{h}, f^{h}\}$ is a Gr\"obner basis for $I^{h}({\bf a})$ with respect to $<_{0}$, so $I^{h}({\bf a})$ is a complete intersection. \hfill $\square$ \par\vspace{0.5\baselineskip}

To illustrate Theorem \ref{Veryfirst}, we present the following example.

\begin{ex1} {\rm Let ${\bf a}=(6851, 5797, 4433, 4862)$. Then $I({\bf a})$ is generated by $x_{2}^{13}-x_1^{11}$, $x_{3}^{17}-x_{1}^{11}$, and $x_{4}^{31}-x_{2}^{13}x_{3}^{17}$. Here, $c_2=d_2=13$ and $c_3=d_3=17$. By Theorem \ref{Veryfirst}, $I^{h}({\bf a})$ is a complete intersection.}
\end{ex1}

Next, we examine the case $d_2<c_2$ and $d_3=c_3$.

\begin{thm1} \label{ACM5} If $d_2<c_{2}$ and $d_{3}=c_3$, then $\overline{C({\bf a})}$ is arithmetically Cohen--Macaulay if and only if $c_{4} \geq c_1+d_{1}+d_{2}$.
\end{thm1}

{\em \noindent Proof.} ($\Leftarrow$) Let $f=x_{4}^{c_4}-x_{1}^{d_1+c_1}x_{2}^{d_2} \in I({\bf a})$. Since $c_{4} \geq c_1+d_{1}+d_{2}$, we have ${\rm in}_{<}(f)=x_{4}^{c_4}$. Then ${\rm in}_{<}(g_1)$, ${\rm in}_{<}(g_2)$, and ${\rm in}_{<}(f)$ are pairwise relatively prime, so $\{g_1, g_2, f\}$ is a Gr\"obner basis for $I({\bf a})$ with respect to $<$ by Remark~\ref{VeryBasicRemark}(1). As $x_1$ divides none of these initial monomials, $\overline{C({\bf a})}$ is arithmetically Cohen--Macaulay by \cite[Theorem 2.2]{HerSta}.\\
($\Rightarrow$) Assume that $c_{4}<c_1+d_{1}+d_{2}$. Since $\overline{C({\bf a})}$ is arithmetically Cohen--Macaulay and ${\rm in}_{<}(f)=x_{1}^{d_1+c_1}x_{2}^{d_2}$, \cite[Theorem~2.2]{HerSta} guarantees the existence of a binomial $B \in I({\bf a})$ such that ${\rm in}_{<}(B)$ divides $x_{2}^{d_2}$. Thus ${\rm in}_{<}(B)=x_{2}^{e}$ for some $e \leq d_2$, which implies $e<c_2$, contradicting the minimality of $c_2$. \hfill $\square$ \par\vspace{0.5\baselineskip}

The following corollary is a consequence of Theorem \ref{ACM5}.

\begin{cor1} \label{Cor12} If $d_2<c_{2}$ and $d_{3}=c_3$, then $I^{h}({\bf a})$ is a complete intersection if and only if $c_{4} \geq c_1+d_{1}+d_{2}$.
\end{cor1}

{\em \noindent Proof.} If $I^{h}({\bf a})$ is a complete intersection, then $\overline{C({\bf a})}$ is arithmetically Cohen--Macaulay, and hence $c_{4} \geq c_1+d_{1}+d_{2}$ by Theorem \ref{ACM5}. Conversely, assume that $c_{4} \geq c_1+d_{1}+d_{2}$. By the proof of Theorem~\ref{ACM5}, $\{g_1,g_2,f\}$ is a Gr\"obner basis for $I({\bf a})$ with respect to $<$. Hence, by Remark~\ref{VeryBasicRemark}(2), $\{g_1^{h},g_2^{h},f^{h}\}$ is a Gr\"obner basis for $I^{h}({\bf a})$ with respect to $<_0$. Consequently, $I^{h}({\bf a})$ is a complete intersection. \hfill $\square$ \par\vspace{0.5\baselineskip}

The following example illustrates Corollary \ref{Cor12}.

\begin{ex1} {\rm Let ${\bf a}=(60, 20, 15, 16)$. Then $I({\bf a})$ is generated by $x_{2}^3-x_1$, $x_{3}^{4}-x_{1}$, and $x_{4}^{5}-x_{2}x_{3}^4$. Here, $d_2=1<3=c_2$ and $d_3=c_3=4$. Since $c_4=5>2=c_1+d_1+d_2$, Corollary~\ref{Cor12} implies that
$I^{h}({\bf a})$ is a complete intersection.}
\end{ex1}

We now consider the case $d_2=c_2$ and $d_3<c_3$.

\begin{thm1} \label{ACM6} If $d_2=c_{2}$ and $d_{3}<c_3$, then $\overline{C({\bf a})}$ is arithmetically Cohen--Macaulay if and only if $c_{4} \geq d_1+c_1+d_{3}$.
\end{thm1}

{\em \noindent Proof.} ($\Leftarrow$) Let $f=x_{4}^{c_4}-x_{1}^{d_1+c_1}x_{3}^{d_3}\in I(\mathbf a)$. Since $c_{4} \geq d_1+c_1+d_{3}$, we have ${\rm in}_{<}(f)=x_{4}^{c_4}$.
Then ${\rm in}_{<}(g_1)$, ${\rm in}_{<}(g_2)$, and ${\rm in}_{<}(f)$ are pairwise relatively prime, so $\{g_1, g_2, f\}$ is a Gr\"obner basis for $I({\bf a})$ with respect to $<$ by Remark~\ref{VeryBasicRemark}(1). As $x_1$ divides none of these initial monomials, $\overline{C({\bf a})}$ is arithmetically Cohen--Macaulay by \cite[Theorem 2.2]{HerSta}.\\
($\Rightarrow$) Assume that $c_{4}<c_1+d_{1}+d_{3}$. Since $\overline{C({\bf a})}$ is arithmetically Cohen--Macaulay and ${\rm in}_{<}(f)=x_{1}^{d_1+c_1}x_{3}^{d_3}$, \cite[Theorem 2.2]{HerSta} guarantees the existence of a binomial $B \in I({\bf a})$ such that ${\rm in}_{<}(B)$ divides $x_{3}^{d_3}$. Thus ${\rm in}_{<}(B)=x_{3}^{e}$ for some $e \leq d_3$, which implies $e<c_3$, contradicting the minimality of $c_3$. \hfill $\square$ \par\vspace{0.5\baselineskip}

Theorem \ref{ACM6} yields the following complete intersection criterion.

\begin{cor1} \label{Cor13} If $d_2=c_{2}$ and $d_{3}<c_3$, then $I^{h}({\bf a})$ is a complete intersection if and only if $c_{4} \geq c_1+d_{1}+d_{3}$.
\end{cor1}

{\em \noindent Proof.} If $I^{h}({\bf a})$ is a complete intersection, then $\overline{C({\bf a})}$ is arithmetically Cohen--Macaulay. By Theorem~\ref{ACM6}, it follows that $c_4 \ge c_1+d_1+d_3$. Conversely, assume that $c_{4} \geq c_1+d_{1}+d_{3}$. By the proof of Theorem~\ref{ACM6}, $\{g_1,g_2,f\}$ is a Gr\"obner basis for $I({\bf a})$ with respect to $<$. Hence, by Remark~\ref{VeryBasicRemark}(2), $\{g_1^{h},g_2^{h},f^{h}\}$ is a Gr\"obner basis for $I^{h}({\bf a})$ with respect to $<_0$. Consequently, $I^{h}({\bf a})$ is a complete intersection. \hfill $\square$ \par\vspace{0.5\baselineskip}

We now study the case $d_2<c_2$ and $d_3<c_3$.

\begin{thm1} \label{Basic4} Suppose that $d_2<c_2$ and $d_{3}<c_3$. Then the ideal $I^{h}({\bf a})$ is a complete intersection if and only if $c_{4} \geq d_{1}+d_{2}+d_{3}$.
\end{thm1}

{\em \noindent Proof.} ($\Leftarrow$) Since $c_{4} \geq d_1+d_2+d_{3}$, we have ${\rm in}_{<}(g_3)=x_{4}^{c_4}$. The initial monomials of $g_1$, $g_2$, and $g_3$ are therefore pairwise
relatively prime. Hence, by Remark~\ref{VeryBasicRemark}(1),
$\{g_i \mid 1 \leq i \leq 3\}$ is a Gr\"obner basis for
$I({\bf a})$ with respect to $<$.
By Remark~\ref{VeryBasicRemark}(2),
$\{g_i^{h} \mid 1 \leq i \leq 3\}$ is a Gr\"obner basis for
$I^{h}({\bf a})$ with respect to $<_0$.
Therefore $I^{h}({\bf a})$ is a complete intersection.\\ 
($\Rightarrow$) Let $T=\{g_{1}^{h}, g_{2}^{h}, x_{4}^{u_4}-x_{0}^{u_0}x_{1}^{u_1}x_{2}^{u_2}x_{3}^{u_3}\}$ be a minimal generating set of $I^{h}({\bf a})$, where $u_{4} \geq u_{1}+u_{2}+u_{3}$. Note that $u_{4} \geq c_{4}$ since $x_{4}^{u_4}-x_{1}^{u_1}x_{2}^{u_2}x_{3}^{u_3} \in I({\bf a})$. Assume that $c_{4}<d_1+d_{2}+d_{3}$. Since $g_{3}^{h}=x_{0}^{d_1+d_{2}+d_{3}-c_{4}}x_{4}^{c_4}-x_{1}^{d_1}x_{2}^{d_{2}}x_{3}^{d_{3}} \in I^{h}({\bf a})$ and the set $T$ generates $I^{h}({\bf a})$, the monomial $x_{4}^{u_4}$ divides $x_{4}^{c_4}$, so $u_4 \leq c_4$, and hence $u_4=c_4$. Moreover, since $d_2<c_2$ and $d_3<c_3$, $x_{1}^{u_1}x_{2}^{u_2}x_{3}^{u_3}$ divides $x_{1}^{d_1}x_{2}^{d_{2}}x_{3}^{d_{3}}$, so $u_i \leq d_i$ for each $1 \leq i \leq 3$. Consequently, $x_{4}^{c_4} (x_1^{d_1-u_1}x_2^{d_2-u_2}x_3^{d_3-u_3}-1) \in I({\bf a})$, which implies that $x_1^{d_1-u_1}x_2^{d_2-u_2}x_3^{d_3-u_3}-1 \in I({\bf a})$, since $I({\bf a})$ is prime and does not contain the variable $x_4$. Therefore, $d_i=u_i$ for all $1 \leq i \leq 3$, and hence $c_4 \geq d_1+d_2+d_3$, a contradiction. \hfill $\square$ \par\vspace{0.5\baselineskip}

When $d_{2}<c_2$, $d_3<c_3$, and $c_4<d_1+d_2+d_3$, the ideal $I^{h}({\bf a})$ is not a complete intersection. Nevertheless, the following Gr\"obner basis computations show that $\overline{C({\bf a})}$ may still be arithmetically Cohen--Macaulay.

Let $w$ be the smallest positive integer such that $c_2-wd_{2} \leq 0$ or $c_3-wd_{3} \leq 0$, namely $w={\rm min} \{l \in \mathbb{N} \setminus \{0\} \mid c_{2}-ld_{2} \leq 0 \ \textrm{or} \ c_{3}-ld_{3} \leq 0\}.$ Since $d_2<c_2$ and $d_3<c_3$, we have $w \geq 2$. For $0 \leq i \leq w-2$ define $$p_{i}=x_{2}^{c_{2}-(i+1)d_{2}}x_{4}^{(i+1)c_{4}}-x_{1}^{c_1}x_{3}^{(i+1)d_{3}}, \ q_{i}=x_{3}^{c_{3}-(i+1)d_3}x_{4}^{(i+1)c_{4}}-x_{1}^{c_1}x_{2}^{(i+1)d_2}.$$ Additionally, define \[
r =
\begin{cases}
x_4^{wc_4}-x_1^{c_1}x_2^{wd_2}x_3^{wd_3-c_3},
& \text{if } c_2-wd_2>0,\\[4pt]
x_4^{wc_4}-x_1^{c_1}x_2^{wd_2-c_2}x_3^{wd_3},
& \text{if } c_2-wd_2 \leq 0.
\end{cases}
\]

\begin{prop1} \label{CMBasic12} Let $d_1=0$, $d_{2}<c_{2}$, $d_{3}<c_{3}$, $c_3>c_2$, and $c_4<d_{2}+d_{3}$. Suppose that the following conditions hold:
\begin{enumerate} \item $(w-1)(c_{4}-d_{2}-d_{3})+c_{2}-c_{1} \geq 0$.
\item $w(c_{4}-d_{2}-d_{3})+c_{3}-c_{1} \geq 0$ when $c_{2}-wd_{2}> 0$.
\item $w(c_{4}-d_{2}-d_{3})+c_{2}-c_{1} \geq 0$ when $c_{2}-wd_{2} \leq 0$.
\end{enumerate}
Then the set \[
\begin{aligned}
T = {} & \{\, g_{1}=x_{2}^{c_2}-x_{1}^{c_1},\;
          g_{2}=x_{3}^{c_3}-x_{1}^{c_1},\;
          f=x_{2}^{d_2}x_{3}^{d_3}-x_{4}^{c_4} \,\} \cup \{\, p_i \mid 0 \le i \le w-2 \,\} \\
      & 
        \cup \{\, q_i \mid 0 \le i \le w-2 \,\}
        \cup \{\, r \,\}.
\end{aligned}
\]
is a Gr\"obner basis for $I({\bf a})$ with respect to $<$.
\end{prop1}

{\em \noindent Proof.} For each $0 \leq i \leq w-2$ the binomials $p_{i}$, $q_{i}$, and $r$ belong to the ideal $I({\bf a})$ by direct computation. Here, ${\rm in}_{<}(g_{1})=x_{2}^{c_{2}}$, ${\rm in}_{<}(g_{2})=x_{3}^{c_{3}}$, and ${\rm in}_{<}(f)=x_{2}^{d_{2}}x_{3}^{d_{3}}$. Since $i+1 \leq w-1$ for all $0 \leq i \leq w-2$, it follows that $(i+1)(c_{4}-d_{2}-d_{3})+c_{2}-c_{1} \geq (w-1)(c_{4}-d_{2}-d_{3})+c_{2}-c_{1}$, as $c_{4}-d_{2}-d_{3}<0$. Hence, $(i+1)(c_{4}-d_{2}-d_{3})+c_{2}-c_{1} \geq 0$, which implies that ${\rm in}_{<}(p_{i})=x_{2}^{c_{2}-(i+1)d_{2}}x_{4}^{(i+1)c_{4}}$. Since $c_3>c_2$, it follows that $(w-1)(c_{4}-d_{2}-d_{3})+c_{3}-c_{1}>(w-1)(c_{4}-d_{2}-d_{3})+c_{2}-c_{1}$. Thus $(w-1)(c_{4}-d_{2}-d_{3})+c_{3}-c_{1} >0$, and therefore ${\rm in}_{<}(q_{i})=x_{3}^{c_{3}-(i+1)d_3}x_{4}^{(i+1)c_{4}}$. Finally, ${\rm in}_{<}(r)=x_{4}^{wc_4}$. Since ${\rm in}_{<}(g_{1})$ and ${\rm in}_{<}(g_{2})$ are relatively prime, we deduce that $S(g_{1},g_{2}) \stackrel{T} {\longrightarrow} 0$. Similarly, we have $S(g_{1},q_{i}) \stackrel{T} {\longrightarrow} 0$ and $S(g_{2},p_i) \stackrel{T} {\longrightarrow} 0$ for all $0 \leq i \leq w-2$, as well as $S(g_{1},r) \stackrel{T} {\longrightarrow} 0$, $S(g_{2},r) \stackrel{T} {\longrightarrow} 0$, and $S(f,r) \stackrel{T} {\longrightarrow} 0$. 

We have that $S(g_{1},f)=p_{0} \stackrel{p_{0}} {\longrightarrow} 0$ and $S(g_{2},f)=q_{0} \stackrel{q_{0}} {\longrightarrow} 0$. We now show that $S(g_{1}, p_{i}) \stackrel{T}{\longrightarrow} 0$ for every $0 \leq i \leq w-2$. Observe that $$S(g_{1},p_{i})=x_{1}^{c_1}(x_{2}^{(i+1)d_2}x_{3}^{(i+1)d_3}-x_{4}^{(i+1)c_4}) \stackrel{f} {\longrightarrow} 0.$$ Similarly, we have $S(g_{2},q_{i})=x_{1}^{c_1}(x_{2}^{(i+1)d_2}x_{3}^{(i+1)d_3}-x_{4}^{(i+1)c_4}) \stackrel{f} {\longrightarrow} 0.$ 

For each $0 \leq i \leq w-3$ we have $S(f,p_{i})=x_{1}^{c_1}x_{3}^{(i+2)d_3}-x_{2}^{c_2-(i+2)d_2}x_{4}^{(i+2)c_4} \stackrel{p_{i+1}} {\longrightarrow} 0$. If $c_{2}-wd_{2} \leq 0$, then $S(f,p_{w-2})=x_{1}^{c_1}x_{2}^{wd_{2}-c_{2}}x_{3}^{wd_3}-x_{4}^{wc_4} \stackrel{r} {\longrightarrow} 0$. On the other hand, if $c_{2}-wd_{2}>0$, then $S(f,p_{w-2})=x_{1}^{c_1}x_{3}^{wd_3}-x_{2}^{c_{2}-wd_2}x_4^{wc_4} \stackrel{r} {\longrightarrow} x_1^{c_1}x_3^{wd_3-c_3}(x_{3}^{c_3}-x_{2}^{c_{2}}) \stackrel{g_1} {\longrightarrow} x_1^{c_1}x_3^{wd_3-c_3}(x_{3}^{c_3}-x_{1}^{c_{1}}) \stackrel{g_2} {\longrightarrow} 0.$

For every $0 \leq i \leq w-3$ we have $S(f,q_{i})=x_{1}^{c_1}x_{2}^{(i+2)d_2}-x_{3}^{c_3-(i+2)d_3}x_{4}^{(i+2)c_4} \stackrel{q_{i+1}} {\longrightarrow} 0$. We now consider the $S$--polynomial $S(f,q_{w-2})$. If $c_{3}-wd_{3}>0$, then \[ S(f,q_{w-2})=x_1^{c_1}x_2^{wd_2}-x_3^{c_3-wd_3}x_4^{wc_4}
\stackrel{r}{\longrightarrow}
x_1^{c_1}x_2^{wd_2-c_2}(x_2^{c_2}-x_3^{c_3})
\stackrel{g_1}{\longrightarrow} \] \[
x_1^{c_1}x_2^{wd_2-c_2}(x_1^{c_1}-x_3^{c_3})
\stackrel{g_2}{\longrightarrow}
0.
\]
Suppose that $c_3-wd_3 \leq 0$. Then $S(f,q_{w-2})=x_1^{c_1}x_2^{wd_2}x_3^{wd_3-c_3}-x_4^{wc_4}$. If $c_2-wd_2>0$, then $S(f,q_{w-2})\stackrel{r}{\longrightarrow}0$. If $c_2-wd_2 \leq 0$, then
\[
S(f,q_{w-2})
\stackrel{r}{\longrightarrow}
x_1^{c_1}x_2^{wd_2-c_2}x_3^{wd_3-c_3}(x_2^{c_2}-x_3^{c_3})
\stackrel{g_1}{\longrightarrow} \] \[
x_1^{c_1}x_2^{wd_2-c_2}x_3^{wd_3-c_3}(x_1^{c_1}-x_3^{c_3})
\stackrel{g_2}{\longrightarrow}
0.\]

Next, we show that $S(p_{i},p_{j}) \stackrel{T}{\longrightarrow} 0$, where $0 \leq i<j \leq w-2$. We compute $S(p_{i},p_{j})=x_{1}^{c_1}x_{3}^{(i+1)d_3}(x_{2}^{(j-i)d_{2}}x_{3}^{(j-i)d_{3}}-x_{4}^{(j-i)c_4}) \stackrel{f}{\longrightarrow} 0$. We now verify that $S(p_i,r) \stackrel{T} {\longrightarrow} 0$ for every $0 \leq i \leq w-2$. Suppose first that $c_2-wd_2>0$. Then $$S(p_i,r)=x_{1}^{c_1}x_{3}^{wd_{3}-c_3}(x_{2}^{c_2+(w-i-1)d_{2}}-x_{3}^{c_{3}-(w-i-1)d_{3}}x_{4}^{(w-i-1)c_{4}}) \stackrel{q_{w-i-2}} {\longrightarrow}$$ $$x_{1}^{c_1}x_{2}^{(w-i-1)d_{2}}x_{3}^{wd_{3}-c_3}(x_{2}^{c_2}-x_{1}^{c_{1}}) \stackrel{g_{1}} {\longrightarrow} 0.$$ Now suppose that $c_2-wd_2 \leq 0$. Then $$S(p_i,r)=x_{1}^{c_1}x_{3}^{(i+1)d_3}(x_{2}^{(w-i-1)d_2}x_{3}^{(w-i-1)d_3}-x_{4}^{(w-i-1)c_4}) \stackrel{f} {\longrightarrow} 0.$$

Next, we show that $S(q_{i},q_{j}) \stackrel{T}{\longrightarrow} 0$, where $0 \leq i<j \leq w-2$. We compute $$S(q_{i},q_{j})=x_{1}^{c_1}x_{2}^{(i+1)d_2}(x_{2}^{(j-i)d_{2}}x_{3}^{(j-i)d_{3}}-x_{4}^{(j-i)c_4}) \stackrel{f}{\longrightarrow} 0.$$ We now verify that $S(q_i,r) \stackrel{T} {\longrightarrow} 0$ for every $0 \leq i \leq w-2$. Suppose first that $c_2-wd_2>0$. Then $S(q_i,r)=x_{1}^{c_1}x_{2}^{(i+1)d_2}(x_{2}^{(w-i-1)d_2}x_{3}^{(w-i-1)d_3}-x_{4}^{(w-i-1)c_4}) \stackrel{f} {\longrightarrow}  0$. Now assume that $c_2-wd_2 \leq 0$, so that 
$r=x_4^{wc_4}-x_1^{c_1}x_2^{wd_2-c_2}x_3^{wd_3}$. Then $$S(q_i,r)=x_{1}^{c_1}x_{2}^{wd_{2}-c_2}(x_{3}^{c_3+(w-i-1)d_{3}}-x_{2}^{c_{2}-(w-i-1)d_{2}}x_{4}^{(w-i-1)c_{4}}) \stackrel{p_{w-i-2}} {\longrightarrow}$$ $$x_{1}^{c_1}x_{2}^{wd_{2}-c_2}x_{3}^{(w-i-1)d_{3}}(x_{3}^{c_3}-x_{1}^{c_{1}}) \stackrel{g_{2}} {\longrightarrow} 0.$$

Finally, we show that $S(p_{i},q_{j}) \stackrel{T}{\longrightarrow} 0$, where $0 \leq i, j \leq w-2$. If $i=j$, then 
$S(p_i,q_i)=x_1^{c_1}(x_2^{c_2}-x_3^{c_3}) \stackrel{g_1} {\longrightarrow} x_1^{c_1}(x_1^{c_1}-x_3^{c_3}) \stackrel{g_2} {\longrightarrow} 0$. Suppose that $i<j$. Then $$S(p_{i},q_{j})=x_{1}^{c_1}(x_{2}^{c_2+(j-i)d_2}-x_{3}^{c_3-(j-i)d_3}x_{4}^{(j-i)c_4}) \stackrel{g_{1}}{\longrightarrow}$$ $$x_{1}^{c_1}(x_{1}^{c_1}x_{2}^{(j-i)d_2}-x_{3}^{c_3-(j-i)d_3}x_{4}^{(j-i)c_4}) \stackrel{q_{j-i-1}}{\longrightarrow} 0.$$ Now suppose that $i>j$. Then $$S(p_{i},q_{j})=x_{1}^{c_1}(x_{2}^{c_2+(j-i)d_2}x_{4}^{(i-j)c_4}-x_{3}^{c_3+(i-j)d_3}) \stackrel{p_{i-j-1}}{\longrightarrow}$$ $$x_{1}^{c_1}x_{3}^{(i-j)d_3}(x_{1}^{c_1}-x_{3}^{c_3}) \stackrel{g_{2}}{\longrightarrow} 0.$$ Thus, $T$ is a Gr\"obner basis for $I({\bf a})$ with respect to $<$. \hfill $\square$ \par\vspace{0.5\baselineskip}


\begin{thm1} \label{CM16} Let $d_1=0$, $d_{2}<c_{2}$, $d_{3}<c_{3}$, $c_3>c_2$, and $c_4<d_{2}+d_{3}$. Suppose that the following conditions hold:
\begin{enumerate} \item $(w-1)(c_{4}-d_{2}-d_{3})+c_{2}-c_{1} \geq 0$.
\item $w(c_{4}-d_{2}-d_{3})+c_{3}-c_{1} \geq 0$ when $c_{2}-wd_{2}> 0$.
\item $w(c_{4}-d_{2}-d_{3})+c_{2}-c_{1} \geq 0$ when $c_{2}-wd_{2} \leq 0$.

\end{enumerate}
Then $\overline{C({\bf a})}$ is arithmetically Cohen--Macaulay.
\end{thm1}

\noindent \textbf{Proof.} By Proposition \ref{CMBasic12}, the set $T$ is a Gr\"obner basis for $I({\bf a})$ with respect to $<$. Since $x_1$ does not divide the initial monomial of any element in $T$, it follows from \cite[Theorem 2.2]{HerSta} that $\overline{C({\bf a})}$ is arithmetically Cohen--Macaulay. \hfill $\square$


\begin{prop1} \label{ACM7} Let $n \geq 3$ be an integer, and let
$a_1=2n(2n+1)$, $a_2=2n+1$, $a_3=2n$, and $a_4=3n+1$. Then
$\overline{C({\bf a})}$ is arithmetically Cohen--Macaulay, and the minimal
number of generators of $I^{h}({\bf a})$ is equal to $2n+2$.
\end{prop1}

{\em \noindent Proof.} Since ${\rm gcd}(a_2,a_3)={\rm gcd}(2n+1,2n)=1$, we obtain
${\rm gcd}(a_1,a_2,a_3,a_4)=1$. The ideal $I({\bf a})$ is generated by
$x_{2}^{2n}-x_1$, $x_{3}^{2n+1}-x_{1}$, and $x_{4}^{2}-x_{2}^{2}x_{3}$. Here $d_2=2<2n=c_2$, $d_3=1<2n+1=c_3$, $c_4=2$,
$w=n$, and $c_2-wd_2=0$. We have
$(w-1)(c_4-d_2-d_3)+c_2-c_1=n>0$ and
$w(c_4-d_2-d_3)+c_2-c_1=n-1>0$. Therefore, by Theorem~\ref{CM16},
$\overline{C({\bf a})}$ is arithmetically Cohen--Macaulay. Moreover, by
Proposition~\ref{CMBasic12}, the set
$$T=\{x_{2}^{2n}-x_1,\; x_{3}^{2n+1}-x_{1},\; x_{2}^{2}x_3-x_{4}^2\}
\cup \{p_{i}=x_{2}^{2n-2(i+1)}x_{4}^{2(i+1)}-x_{1}x_3^{i+1}
\mid 0 \leq i \leq n-2\}$$
$$\cup \{q_{i}=x_{3}^{2n+1-(i+1)}x_{4}^{2(i+1)}-x_{1}x_{2}^{2(i+1)}
\mid 0 \leq i \leq n-2\}
\cup \{r=x_{4}^{2n}-x_{1}x_{3}^{n}\}$$
is a Gr\"obner basis for $I({\bf a})$ with respect to $<$. By Remark~\ref{VeryBasicRemark} (2), the homogenizations of the elements of $T$
form a Gr\"obner basis of $I^{h}({\bf a})$ with respect to $<_0$, namely
$$\{x_{2}^{2n}-x_{0}^{2n-1}x_1,\; x_{3}^{2n+1}-x_{0}^{2n}x_1,\;
x_{2}^{2}x_3-x_{0}x_{4}^{2}\}
\cup \{p^{h}_{i} \mid 0 \leq i \leq n-2\}$$
$$\cup \{q^{h}_{i} \mid 0 \leq i \leq n-2\} \cup \{r^{h}\}.$$ A direct computation shows that the above set is a minimal
generating set of $I^{h}({\bf a})$. Thus $\mu(I^{h}({\bf a}))=2n+2$.
\hfill $\square$  \par\vspace{0.5\baselineskip}

{\em Finally, we consider the case $a_{4}={\rm max}\{a_1,\ldots,a_4\}$. Without loss of generality, we fix the assumptions $c_{3}>c_{1}$, $c_{2}>c_{1}$, $d_{2} \leq c_2$, and $d_{3} \leq c_3$.} Under these assumptions, we give a characterization for when $I^{h}({\bf a})$ is a complete intersection.

\begin{thm1} \label{Comp2} The ideal $I^{h}({\bf a})$ is a complete intersection if and only if $d_{2}=c_2$ and $d_{3}=c_{3}$.
\end{thm1}

{\em \noindent Proof.}  ($\Leftarrow$) Since $a_{4}={\rm max}\{a_1,\ldots,a_4\}$, we use the degree reverse lexicographic order $<$ with $x_{1}>x_2>x_3>x_4$. Let $f=x_{1}^{d_1+2c_1}-x_{4}^{c_4} \in I({\bf a})$. Then ${\rm in}_{<}(g_1)=x_{2}^{c_2}$, ${\rm in}_{<}(g_2)=x_{3}^{c_3}$, and ${\rm in}_{<}(f)=x_{1}^{d_{1}+2c_1}$ are pairwise relatively prime. Hence, by Remark \ref{VeryBasicRemark}(1), the set $\{g_1, g_2, f\}$ is a Gr\"obner basis for $I({\bf a})$ with respect to $<$. By Remark \ref{VeryBasicRemark}(2), the set $\{g_{1}^{h}, g_{2}^{h}, f^{h}\}$ is a Gr\"obner basis for $I^{h}({\bf a})$ with respect to $<_{0}$. Consequently, $I^{h}({\bf a})$ is a complete intersection.\\
($\Rightarrow$) Let $T=\{g_{1}^{h}, g_{2}^{h}, x_{1}^{u_1}-x_{0}^{u_{0}}x_{2}^{u_2}x_{3}^{u_3}x_{4}^{u_4}\}$ be a minimal generating set of $I^{h}({\bf a})$, where $u_{1} \geq c_{1}$ and $u_{1} \geq u_2+u_3+u_4$. Since $a_{4}={\rm max}\{a_1,\ldots,a_4\}$, we have $c_{4}<d_{1}+d_{2}+d_{3}$, and therefore $x_{0}^{d_{1}+d_{2}+d_{3}-c_{4}}x_{4}^{c_{4}}-x_{1}^{d_{1}}x_{2}^{d_{2}}x_{3}^{d_{3}} \in I^{h}({\bf a})$. Since the set $T$ generates $I^{h}({\bf a})$, the monomial $x_{0}^{u_{0}}x_{2}^{u_2}x_{3}^{u_3}x_{4}^{u_4}$ divides $x_{0}^{d_{1}+d_{2}+d_{3}-c_{4}}x_{4}^{c_{4}}$. Thus, $u_{2}=u_{3}=0$ and $u_{4} \leq c_{4}$. Hence, $x_{1}^{u_1}-x_{0}^{u_{0}}x_{4}^{u_4} \in I^{h}({\bf a})$, so $x_{1}^{u_1}-x_{4}^{u_4} \in I({\bf a})$, and therefore $u_{4} \geq c_{4}$. This implies $u_{4}=c_{4}$. Since at least one of $d_2$ and $d_3$ is nonzero, we may assume without loss of generality that $d_{2} \neq 0$. Observe that $x_{1}^{c_1+d_1}x_{3}^{d_3}-x_{2}^{c_{2}-d_2}x_{1}^{u_1} \in I({\bf a})$. As $c_2-d_2<c_2$, we conclude that $u_{1} \geq c_1+d_{1}$, and therefore $x_{3}^{d_3}-x_{1}^{u_1-c_1-d_1}x_{2}^{c_2-d_2} \in I({\bf a})$. Thus $d_{3} \geq c_3$, and since by assumption $d_3 \leq c_3$, we conclude that $d_{3}=c_3$. Now consider the binomial $x_{1}^{c_1+d_1}x_{2}^{d_2}-x_{3}^{c_3-d_3}x_{4}^{c_4} \in I({\bf a})$. Since $c_3=d_3$ and $x_{1}^{u_1}-x_{4}^{c_4} \in I({\bf a})$, we have $x_{1}^{c_1+d_1}x_{2}^{d_2}-x_{1}^{u_1} \in I({\bf a})$, which implies $x_{2}^{d_2}-x_{1}^{u_1-c_1-d_1} \in I({\bf a})$. Therefore, $d_2 \geq c_2$. Since $d_2 \leq c_2$ by assumption, we conclude that $d_2=c_2$. \hfill $\square$ \par\vspace{0.5\baselineskip}

The next example demonstrates the complete intersection property of $I^h(\mathbf{a})$ when $d_2=c_2$ and $d_3=c_3$.

\begin{ex1} {\rm Let ${\bf a}=(140, 84, 60, 245)$. Then $I({\bf a})$ is generated by the binomials $x_{2}^{5}-x_{1}^{3}$, $x_{3}^{7}-x_{1}^3$, and $x_{4}^4-x_1x_2^5x_3^7$. Here, $c_{2}=d_{2}=5$ and $c_{3}=d_{3}=7$. By Theorem \ref{Comp2}, the ideal $I^{h}({\bf a})$ is a complete intersection.}
\end{ex1}



\section{Projective closures in case 3}

In this section, we study Case 3 of Theorem \ref{CasesMinimal}, where $I({\bf a})$ has a minimal generating set $\{x_{2}^{c_2}-x_{1}^{c_1}, x_{3}^{c_3}-x_{1}^{b_1}x_{2}^{b_2}, x_{4}^{c_4}-x_{1}^{d_1}x_{2}^{d_2}x_{3}^{d_3}\}$. Both $b_1$ and $b_2$ are nonzero, and at least two of $d_1$, $d_2$, and $d_3$ are nonzero.

Suppose first that $a_{i}={\rm max}\{a_1,\ldots,a_4\}$ for some $i \in \{1,2\}$. {\em Without loss of generality, we assume $a_{1}={\rm max}\{a_1,\ldots,a_4\}$ and $b_2 \leq c_2$.} The next theorem provides a complete intersection criterion for $I^{h}({\bf a})$ when $b_2<c_2$.

\begin{thm1} \label{Basic5} Suppose that $b_2<c_2$. Then $I^{h}({\bf a})$ is a complete intersection if and only if $c_{3} \geq b_{1}+b_{2}$ and the element $c_{4}a_{4} \in \mathbb{N}\{a_1,\ldots,a_4\}$ admits a factorization $c_{4}a_{4}=v_{1}a_{1}+v_{2}a_{2}+v_{3}a_{3}$ with $c_{4} \geq v_1+v_{2}+v_{3}$, for some $v_1, v_2, v_3 \in \mathbb{N}$. 
\end{thm1}

{\em \noindent Proof.} ($\Leftarrow$) Let $g_{1}=x_{2}^{c_2}-x_{1}^{c_1}, g_{2}=x_{3}^{c_3}-x_{1}^{b_1}x_{2}^{b_2}$, and $g_{3}=x_{4}^{c_4}-x_{1}^{v_1}x_{2}^{v_2}x_{3}^{v_3} \in I({\bf a})$, where $c_{4} \geq v_1+v_{2}+v_{3}$. Consider the degree reverse lexicographic order $<$ with $x_{4}>x_{3}>x_{2}>x_{1}$. Then ${\rm in}_{<}(g_1)=x_{2}^{c_2}$, ${\rm in}_{<}(g_2)=x_{3}^{c_3}$, and ${\rm in}_{<}(g_{3})=x_{4}^{c_4}$ are pairwise relatively prime. Hence, by Remark \ref{VeryBasicRemark}(1), $\{g_{i} \mid 1 \leq i \leq 3\}$ is a Gr\"obner basis for $I({\bf a})$ with respect to $<$. By Remark \ref{VeryBasicRemark}(2), $\{g_{i}^{h} \mid 1 \leq i \leq 3\}$ is a Gr\"obner basis for $I^{h}({\bf a})$ with respect to $<_{0}$. Therefore $I^{h}({\bf a})$ is a complete intersection.\\
($\Rightarrow$) Assume first that $c_3<b_1+b_2$. Since $\overline{C({\bf a})}$ is arithmetically Cohen--Macaulay and ${\rm in}_{<}(g_2)=x_1^{b_1}x_2^{b_2}$, \cite[Theorem~2.2]{HerSta} yields a binomial $B\in I({\bf a})$
whose initial monomial divides $x_2^{b_2}$, contradicting
the minimality of $c_2$. Now suppose that for every factorization $c_{4}a_{4}=v_{1}a_{1}+v_{2}a_{2}+v_{3}a_{3}$ we have $c_{4}<v_1+v_{2}+v_{3}$. Let $T=\{g_{1}^{h}, g_{2}^{h}, x_{4}^{z_4}-x_{0}^{z_{0}}x_{1}^{z_1}x_{2}^{z_2}x_{3}^{z_3}\}$ be a minimal generating set of $I^{h}({\bf a})$, where $z_{4} \geq z_{1}+z_{2}+z_{3}$. Since $x_{4}^{z_4}-x_{1}^{z_1}x_{2}^{z_2}x_{3}^{z_3} \in I({\bf a})$, we have $z_{4} \geq c_4$, and $c_4<z_1+z_2+z_3$ implies $z_4>c_4$. Moreover, $x_{0}^{d_1+d_2+d_3-c_4}x_{4}^{c_4}-x_{1}^{d_1}x_2^{d_2}x_{3}^{d_3} \in I^{h}({\bf a})$, and since $T$ generates $I^{h}({\bf a})$, we have $x_{4}^{z_4}$ divides $x_{4}^{c_4}$, hence $z_{4} \leq c_4$, a contradiction. \hfill $\square$ \par\vspace{0.5\baselineskip}

The failure of either condition in Theorem \ref{Basic5} prevents $I^{h}({\bf a})$ from being a complete intersection. The following example illustrates this situation.

\begin{ex1} {\rm Let ${\bf a}=(a_1=12, a_2=4, a_3=10, a_4=11)$. Then $I({\bf a})$ is generated by $x_{2}^{3}-x_1$, $x_{3}^{2}-x_{1}x_{2}^{2}$, and $x_{4}^{2}-x_{2}^{3}x_{3}$. Here, $b_{2}=2<3=c_{2}$, $c_{3}=2<3=b_{1}+b_{2}$, and $2a_4=a_1+a_3$. By Theorem \ref{Basic5}, $I^{h}({\bf a})$ is not a complete intersection.}
\end{ex1}


Next, we provide a complete intersection criterion for $I^{h}({\bf a})$ when $b_2=c_2$.

\begin{thm1} \label{Basic23} Suppose that $b_2=c_2$. Then $I^{h}({\bf a})$ is a complete intersection if and only if the element $c_{4}a_{4} \in \mathbb{N}\{a_1,\ldots,a_4\}$ admits a factorization $c_{4}a_{4}=v_{1}a_{1}+v_{2}a_{2}+v_{3}a_{3}$ for some $v_1, v_2, v_3 \in \mathbb{N}$ with $c_{4} \geq v_1+v_{2}+v_{3}$. 
\end{thm1}


{\em \noindent Proof.} ($\Leftarrow$) Let $g_{1}=x_{2}^{c_2}-x_{1}^{c_1}, g_{2}=x_{3}^{c_3}-x_{1}^{b_1+c_1}$, and $g_{3}=x_{4}^{c_4}-x_{1}^{v_1}x_{2}^{v_2}x_{3}^{v_3} \in I({\bf a})$, where $c_{4} \geq v_1+v_{2}+v_{3}$. Consider the degree reverse lexicographic order $<$ with $x_{4}>x_{3}>x_{2}>x_{1}$. Then, ${\rm in}_{<}(g_1)=x_{2}^{c_2}$, ${\rm in}_{<}(g_2)=x_{3}^{c_3}$, and ${\rm in}_{<}(g_3)=x_{4}^{c_4}$ are pairwise relatively prime. Hence, by Remark \ref{VeryBasicRemark}(1), $\{g_{i} \mid 1 \leq i \leq 3\}$ is a Gr\"obner basis for $I({\bf a})$ with respect to $<$. By Remark \ref{VeryBasicRemark}(2), $\{g_{i}^{h} \mid 1 \leq i \leq 3\}$ is a Gr\"obner basis for $I^{h}({\bf a})$ with respect to $<_{0}$. Thus, $I^{h}({\bf a})$ is a complete intersection.\\ ($\Rightarrow$) The proof of this direction follows the same argument as the $(\Rightarrow)$ direction of Theorem~\ref{Basic5}, using only the obstruction arising from the factorization of $c_4a_4$, and is therefore omitted. \hfill $\square$ \par\vspace{0.5\baselineskip}

The following example illustrates Theorem~\ref{Basic23}.


\begin{ex1} {\rm Let ${\bf a}=(a_{1}=18, a_{2}=9, a_{3}=12, a_{4}=14)$. Then $I({\bf a})$ is generated by $x_{2}^{2}-x_{1}$, $x_{3}^{3}-x_{1}x_{2}^{2}$, and $x_{4}^{3}-x_{2}^2x_{3}^{2}$. Here, $b_{2}=c_{2}=2$ and $3a_4=a_1+2a_3$. By Theorem \ref{Basic23}, $I^{h}({\bf a})$ is a complete intersection.} 
\end{ex1}

{\em Suppose now that $a_{3}={\rm max}\{a_1,\ldots,a_4\}$. Without loss of generality, we fix the assumption $c_{2}>c_{1}$.} The following theorem establishes a complete intersection criterion in the case where $c_4 \geq d_1+d_2+d_3$.

\begin{thm1} \label{basicvfff} Suppose that $c_{4} \geq d_1+d_2+d_3$. Then $I^{h}({\bf a})$ is a complete intersection if and only if the element $c_{3}a_3 \in \mathbb{N}\{a_1,\ldots,a_4\}$ admits a factorization $c_{3}a_3=u_1a_{1}$ for some $u_1 \in \mathbb{N}$ with $u_{1}>c_1$.
\end{thm1}
{\em \noindent Proof.} ($\Leftarrow$) Let $f_{1}=x_{2}^{c_2}-x_{1}^{c_1}, f_{2}=x_{4}^{c_4}-x_{1}^{d_1}x_2^{d_2}x_{3}^{d_3}, f_{3}=x_{1}^{u_1}-x_{3}^{c_3} \in I({\bf a})$. Consider the degree reverse lexicographic order $<$ with $x_4>x_2>x_1>x_3$. Since $c_{4} \geq d_1+d_2+d_3$, we have ${\rm in}_{<}(f_{2})=x_{4}^{c_4}$. Then ${\rm in}_{<}(f_{1})=x_{2}^{c_2}$, ${\rm in}_{<}(f_{2})$, and ${\rm in}_{<}(f_{3})=x_{1}^{u_1}$ are pairwise relatively prime. By Remark \ref{VeryBasicRemark}(1), $\{f_i \mid 1 \leq i \leq 3\}$ is a Gr\"obner basis for $I({\bf a})$ with respect to $<$. By Remark \ref{VeryBasicRemark}(2), $\{f_{i}^h \mid 1 \leq i \leq 3\}$ is a Gr\"obner basis for $I^{h}({\bf a})$ with respect to $<_{0}$. Thus, $I^{h}({\bf a})$ is a complete intersection.\\
($\Rightarrow$) Let $T=\{x_{1}^{u_1}-x_{0}^{u_0}x_{2}^{u_2}x_{3}^{u_3}x_{4}^{u_4}, x_{2}^{c_2}-x_{0}^{c_2-c_1}x_{1}^{c_1}, x_{4}^{c_4}-x_{0}^{c_4-d_1-d_2-d_3}x_{1}^{d_1}x_{2}^{d_2}x_{3}^{d_3}\}$ be a minimal generating set of $I^{h}({\bf a})$, where $u_{1} \geq u_2+u_3+u_4$. Since $a_{3}={\rm max}\{a_1,\ldots,a_4\}$, we have $c_{3}<b_{1}+b_{2}$, and therefore $x_{0}^{b_{1}+b_{2}-c_{3}}x_{3}^{c_3}-x_{1}^{b_1}x_{2}^{b_2} \in I^{h}({\bf a})$. Furthermore, since $T$ generates $I^{h}({\bf a})$, at least one of $x_{0}^{u_0}x_{2}^{u_2}x_{3}^{u_3}x_{4}^{u_4}$ and $x_{0}^{c_4-d_1-d_2-d_3}x_{1}^{d_1}x_{2}^{d_2}x_{3}^{d_3}$ divides $x_{0}^{b_{1}+b_{2}-c_{3}}x_{3}^{c_3}$. We distinguish two cases.\\
(1) Suppose that $x_{0}^{u_0}x_{2}^{u_2}x_{3}^{u_3}x_{4}^{u_4}$ divides $x_{0}^{b_{1}+b_{2}-c_{3}}x_{3}^{c_3}$. Then $u_{2}=u_{4}=0$ and $u_{3} \leq c_{3}$, so $x_{1}^{u_1}-x_{0}^{u_0}x_{3}^{u_3} \in I^{h}({\bf a})$, implying $x_{1}^{u_1}-x_{3}^{u_3} \in I({\bf a})$. Thus $u_{3} \geq c_{3}$ by the minimality of $c_3$, hence $u_{3}=c_{3}$. Since $c_{1}a_1 \neq c_{3}a_3$, we conclude that $u_1>c_1$.\\
(2) Suppose that $x_{0}^{c_4-d_1-d_2-d_3}x_{1}^{d_1}x_{2}^{d_2}x_{3}^{d_3}$ divides $x_{0}^{b_{1}+b_{2}-c_{3}}x_{3}^{c_3}$. Since the latter monomial contains neither $x_1$ nor $x_2$, it follows that $d_1=d_2=0$, contradicting the assumption that at least two among $d_1$, $d_2$, $d_3$ are strictly positive.  \hfill $\square$ \par\vspace{0.5\baselineskip}

The following example illustrates Theorem~\ref{basicvfff}.

\begin{ex1} {\rm Let ${\bf a}=(a_{1}=4845,a_2=4080,a_3=5491,a_{4}=4883)$. Then $I({\bf a})$ is generated by $x_{2}^{19}-x_{1}^{16}$, $x_{3}^{15}-x_{1}x_{2}^{19}$, and $x_{4}^{17}-x_{1}^{16}x_{3}$. Here $c_4=17=d_1+d_2+d_3$ and $15a_3=17a_1 $. By Theorem \ref{basicvfff}, $I^{h}({\bf a})$ is a complete intersection.}
\end{ex1}

When $c_4<d_1+d_2+d_3$, the complete intersection property of the projective closure depends on the factorizations of both $c_3a_3$ and $c_4a_4$.

\begin{thm1} \label{Basic24} Suppose that $c_4<d_1+d_2+d_3$. Then $I^{h}({\bf a})$ is a complete intersection if and only if at least one of the following conditions holds: \begin{enumerate} \item[(A)] The element $c_{3}a_3 \in \mathbb{N}\{a_1,\ldots,a_4\}$ admits a factorization $c_{3}a_3=u_1a_{1}$ for some integer $u_{1}>c_1$, and the element $c_{4}a_4 \in \mathbb{N}\{a_1,\ldots,a_4\}$ admits a factorization $c_{4}a_4=v_1a_{1}+v_2a_{2}+v_3a_3$ for some $v_1, v_2, v_3 \in \mathbb{N}$ with $c_4 \geq v_1+v_2+v_3$.
\item[(B)] The element $c_{3}a_3 \in \mathbb{N}\{a_1,\ldots,a_4\}$ admits a factorization $c_{3}a_3=v_4a_{4}$ for some integer $v_{4}>c_4$, and the element $c_{4}a_4 \in \mathbb{N}\{a_1,\ldots,a_4\}$ admits a factorization $c_{4}a_4=u_1a_{1}$ for some $u_1 \in \mathbb{N}$ such that $u_1>c_4$ and $u_1>c_1$.
\end{enumerate}
\end{thm1}
{\em \noindent Proof.} $(\Leftarrow)$ Suppose first that condition (A) holds. Let $f_{1}=x_{2}^{c_2}-x_{1}^{c_1}, f_{2}=x_{1}^{u_1}-x_{3}^{c_3} \in I({\bf a})$, and $f_{3}=x_{4}^{c_4}-x_{1}^{v_1}x_{2}^{v_2}x_{3}^{v_3} \in I({\bf a})$, where $c_{4} \geq v_1+v_{2}+v_{3}$. Consider the degree reverse lexicographic order $<$ with $x_4>x_2>x_1>x_3$. Then ${\rm in}_{<}(f_{1})=x_{2}^{c_2}$, ${\rm in}_{<}(f_{2})=x_{1}^{u_1}$, and ${\rm in}_{<}(f_{3})=x_{4}^{c_4}$ are pairwise relatively prime. Hence, by Remark \ref{VeryBasicRemark}(1), $\{f_{i} \mid 1 \leq i \leq 3\}$ is a Gr\"obner basis for $I({\bf a})$ with respect to $<$. By Remark \ref{VeryBasicRemark}(2), $\{f_1^h,f_2^h,f_3^h\}$ is a Gr\"obner basis for $I^{h}({\bf a})$ with respect to $<_{0}$. Thus, $I^{h}({\bf a})$ is a complete intersection. 

Now suppose that condition (B) holds. Consider the binomials $f_{4}=x_{4}^{v_4}-x_{3}^{c_3} \in I({\bf a})$ and $f_{5}=x_{1}^{u_1}-x_{4}^{c_4} \in I({\bf a})$, where $u_1>c_4$. Then ${\rm in}_{<}(f_{1})$, ${\rm in}_{<}(f_{4})=x_{4}^{v_4}$, and ${\rm in}_{<}(f_{5})=x_{1}^{u_1}$ are pairwise relatively prime. By Remark \ref{VeryBasicRemark}(1), $\{f_1, f_4, f_5\}$ is a Gr\"obner basis for $I({\bf a})$ with respect to $<$. By Remark \ref{VeryBasicRemark}(2), $\{f_1^h, f_4^h, f_5^h\}$ is a Gr\"obner basis for $I^{h}({\bf a})$ with respect to $<_{0}$. Consequently, $I^{h}({\bf a})$ is a complete intersection.\\
$(\Rightarrow)$ Let $T=\{x_{1}^{u_1}-x_{0}^{u_0}x_{2}^{u_2}x_{3}^{u_3}x_{4}^{u_4}, x_{2}^{c_2}-x_{0}^{c_2-c_1}x_{1}^{c_1}, x_{4}^{v_4}-x_{0}^{v_0}x_{1}^{v_1}x_{2}^{v_2}x_{3}^{v_3}\}$ be a minimal generating set of $I^{h}({\bf a})$, where $u_{1} \geq u_2+u_3+u_4$ and $v_{4} \geq v_1+v_2+v_3$. Since $c_{3}<b_{1}+b_{2}$, the binomial $x_{0}^{b_{1}+b_{2}-c_{3}}x_{3}^{c_3}-x_{1}^{b_1}x_{2}^{b_2}$ belongs to $I^{h}({\bf a})$. Since $T$ generates $I^{h}({\bf a})$, at least one of $x_{0}^{u_0}x_{2}^{u_2}x_{3}^{u_3}x_{4}^{u_4}$ and $x_{0}^{v_0}x_{1}^{v_1}x_{2}^{v_2}x_{3}^{v_3}$ divides $x_{0}^{b_{1}+b_{2}-c_{3}}x_{3}^{c_3}$. We distinguish two cases.\\
(1) Suppose that $x_{0}^{u_0}x_{2}^{u_2}x_{3}^{u_3}x_{4}^{u_4}$ divides $x_{0}^{b_{1}+b_{2}-c_{3}}x_{3}^{c_3}$. Then $u_{2}=u_{4}=0$ and $u_{3} \leq c_{3}$. Hence, $x_{1}^{u_1}-x_{0}^{u_0}x_{3}^{u_3} \in I^{h}({\bf a})$, and therefore $x_{1}^{u_1}-x_{3}^{u_3} \in I({\bf a})$. Thus $u_{3} \geq c_{3}$ and therefore $u_{3}=c_{3}$. Since $c_1a_1 \neq c_3a_3$, we have $u_1>c_1$. Furthermore, the binomial $x_{0}^{d_1+d_2+d_3-c_4}x_{4}^{c_4}-x_{1}^{d_1}x_2^{d_2}x_3^{d_3}$ belongs to $I^{h}({\bf a})$. Since $T$ generates $I^{h}({\bf a})$, the monomial $x_4^{v_4}$ must divide $x_{0}^{d_1+d_2+d_3-c_4}x_{4}^{c_4}$, and therefore $v_4 \leq c_4$. But $x_{4}^{v_4}-x_{1}^{v_1}x_{2}^{v_2}x_{3}^{v_3} \in I({\bf a})$, so $v_{4} \geq c_4$, and thus $v_4=c_4$. Consequently, the condition (A) holds.\\
(2) Suppose that $x_{0}^{v_0}x_{1}^{v_1}x_{2}^{v_2}x_{3}^{v_3}$ divides $x_{0}^{b_{1}+b_{2}-c_{3}}x_{3}^{c_3}$. Then $v_{1}=v_{2}=0$ and $v_{3} \leq c_{3}$, so $x_{4}^{v_4}-x_{0}^{v_0}x_{3}^{v_3} \in I^{h}({\bf a})$, implying $x_{4}^{v_4}-x_{3}^{v_3} \in I({\bf a})$. Thus $v_{3} \geq c_3$, and therefore $v_{3}=c_{3}$. Since $c_{3}a_{3} \neq c_4a_4$, it follows that $v_4>c_{4}$. Next, consider the binomial $x_{0}^{d_1+d_2+d_3-c_4}x_{4}^{c_4}-x_{1}^{d_1}x_2^{d_2}x_3^{d_3} \in I^{h}({\bf a})$. Since $T$ generates $I^{h}({\bf a})$, the monomial $x_{0}^{u_0}x_{2}^{u_2}x_{3}^{u_3}x_{4}^{u_4}$ must divide $x_{0}^{d_1+d_2+d_3-c_4}x_{4}^{c_4}$, so $u_2=u_3=0$ and $u_4 \leq c_4$. Hence, $x_{1}^{u_1}-x_{0}^{u_0}x_{4}^{u_4} \in I^{h}({\bf a})$ with $u_1 \geq u_4$, implying $x_{1}^{u_1}-x_{4}^{u_4} \in I({\bf a})$ with $u_4 \geq c_4$. Thus $u_4=c_4$. So $u_1 \geq c_4$. Since $a_1 \neq a_4$, it follows that $u_1>c_4$, and because $c_1a_1 \neq c_4a_4$, we also have $u_1>c_1$. \begin{rem1}\label{rem:Basic24}
{\rm Conditions {\rm (A)} and {\rm (B)} in Theorem~\ref{Basic24} can hold simultaneously. Consider
${\bf a}=(a_{1}=3, a_{2}=12, a_{3}=48, a_{4}=8)$. Then $
I({\bf a})$ is generated by $x_2-x_1^4, x_3-x_1^4x_2^3,$ and $x_4^3-x_1^4x_2.$ Hence
$ c_1=4, c_2=c_3=1, c_4=3, b_1=4,b_2=3, d_1=4, d_2=1, d_3=0,
$ and $c_4<d_1+d_2+d_3.$ Condition {\rm (A)} holds because
$c_3a_3=48=16a_1$ and $c_4a_4=24=2a_2.$ Condition {\rm (B)} holds because
$c_3a_3=48=6a_4$ and $c_4a_4=24=8a_1$.}
\end{rem1}

The following examples illustrate some of the cases arising in Theorem~\ref{Basic24}.

\begin{ex1} {\rm Let ${\bf a}=(a_{1}=5083,a_2=4641,a_3=8602,a_{4}=8372)$. Then $I({\bf a})$ is generated by $x_{2}^{23}-x_{1}^{21}$, $x_{3}^{13}-x_{1}x_{2}^{23}$, and $x_{4}^{17}-x_{1}^{7}x_{2}^{23}$. In this case, we have $c_1=21$, $c_3=13$, and $c_4=17<30=d_1+d_2+d_3$. Note that there are no $v_1, v_2, v_3 \in \mathbb{N}$ such that $17a_{4}=v_1a_{1}+v_{2}a_{2}+v_{3}a_{3}$ with $v_1+v_2+v_3 \leq 17,$ and thus condition (A) of Theorem~\ref{Basic24} does not hold. Additionally, although $17a_4=28a_1$, the element $13a_3$ does not admit a factorization $13a_3=v_4a_4$ with $v_4>17$, and therefore condition (B) also fails. By Theorem \ref{basicvfff}, $I^{h}({\bf a})$ is not a complete intersection.}
\end{ex1}

\begin{ex1} {\rm Let ${\bf a}=(a_{1}=7429,a_2=6783,a_3=21850,a_{4}=9775)$. Then $I({\bf a})$ is generated by $x_{2}^{23}-x_{1}^{21}$, $x_{3}^{17}-x_{1}^{29}x_{2}^{23}$, and $x_{4}^{19}-x_{1}^{4}x_{2}^{23}$. In this case, we have $c_1=21$, $c_2=23$, $c_3=17$, and $c_4=19<27=d_1+d_2+d_3$. Note that there do not exist $v_1, v_2, v_3 \in \mathbb{N}$ such that $19a_{4}=v_1a_{1}+v_{2}a_{2}+v_{3}a_{3}$ with $v_1+v_2+v_3 \leq 19,$ so condition (A) of Theorem~\ref{Basic24} does not hold. On the other hand, observe that $17a_3=38a_4$ and $19a_4=25a_1$. Therefore, condition (B) of Theorem~\ref{Basic24} is satisfied, and it follows that $I^{h}({\bf a})$ is a complete intersection.}
\end{ex1}

{\em Suppose now that $a_{4}={\rm max}\{a_1,\ldots,a_4\}$. Without loss of generality, we fix $c_{2}>c_{1}$.} The next theorem gives a complete intersection criterion for $c_{3} \geq b_1+b_2$.

\begin{thm1} \label{Basic25} Suppose that $c_{3} \geq b_1+b_2$. Then $I^{h}({\bf a})$ is a complete intersection if and only if the element $c_{4}a_4 \in \mathbb{N}\{a_1,\ldots,a_4\}$ admits a factorization $c_{4}a_4=u_1a_{1}$ for some $u_1 \in \mathbb{N}$ with $u_{1}>c_1$.
\end{thm1}

{\em \noindent Proof.} ($\Leftarrow$) Let $f_{1}=x_{2}^{c_2}-x_{1}^{c_1}, f_{2}=x_{3}^{c_3}-x_{1}^{b_1}x_2^{b_2}$, and $f_{3}=x_{1}^{u_1}-x_{4}^{c_4} \in I({\bf a})$. Consider the degree reverse lexicographic order $<$ with $x_1>x_2>x_3>x_4$. Since $c_3 \geq b_1+b_2$, we have ${\rm in}_{<}(f_2)=x_{3}^{c_3}$. Then ${\rm in}_{<}(f_{1})=x_{2}^{c_2}$, ${\rm in}_{<}(f_{2})$, and ${\rm in}_{<}(f_{3})=x_{1}^{u_1}$ are pairwise relatively prime. Hence, by Remark \ref{VeryBasicRemark}(1), $\{f_{i} \mid 1 \leq i \leq 3\}$ is a Gr\"obner basis for $I({\bf a})$ with respect to $<$. By Remark \ref{VeryBasicRemark}(2), $\{f_{i}^h \mid 1 \leq i \leq 3\}$ is a Gr\"obner basis for $I^{h}({\bf a})$ with respect to $<_{0}$. Hence, $I^{h}({\bf a})$ is a complete intersection.

($\Rightarrow$) Let $T=\{x_{1}^{u_1}-x_{0}^{u_0}x_{2}^{u_2}x_{3}^{u_3}x_{4}^{u_4}, x_{2}^{c_2}-x_{0}^{c_2-c_1}x_{1}^{c_1}, x_{3}^{c_3}-x_{0}^{c_3-b_1-b_2}x_{1}^{b_1}x_{2}^{b_2}\}$ be a minimal generating set of $I^{h}({\bf a})$, where $u_{1} \geq u_2+u_3+u_4$. Since $a_{4}={\rm max}\{a_1,\ldots,a_4\}$, it follows that $c_{4}<d_{1}+d_{2}+d_3$, so $x_{0}^{d_{1}+d_{2}+d_3-c_{4}}x_{4}^{c_4}-x_{1}^{d_1}x_{2}^{d_2}x_{3}^{d_3} \in I^{h}({\bf a})$. Because $T$ generates $I^{h}({\bf a})$, the monomial $x_{0}^{u_0}x_{2}^{u_2}x_{3}^{u_3}x_{4}^{u_4}$ must divide $x_{0}^{d_{1}+d_{2}+d_3-c_{4}}x_{4}^{c_4}$. Hence, $u_{2}=u_{3}=0$ and $u_{4} \leq c_4$, so $x_{1}^{u_1}-x_{0}^{u_0}x_{4}^{u_4} \in I^{h}({\bf a})$, and therefore $x_{1}^{u_1}-x_{4}^{u_4} \in I({\bf a})$. Thus $u_{4} \geq c_{4}$, and hence $u_{4}=c_{4}$. Since $c_{1}a_1 \neq c_{4}a_4$, we conclude that $u_1>c_1$. \hfill $\square$ \par\vspace{0.5\baselineskip}

The following example illustrates Theorem \ref{Basic25}.

\begin{ex1} {\rm Let ${\bf a}=(a_{1}=455,a_2=385,a_3=120,a_{4}=1092)$. Then $I({\bf a})$ is generated by $x_{2}^{13}-x_{1}^{11}$, $x_{3}^{7}-x_{1}x_{2}$, and $x_{4}^{5}-x_{1}x_2^{13}$. Here, $c_1=11$, $c_{3}=7> 2=b_1+b_2$, and $c_{4}=5$. Also $5a_4=12a_1$. By Theorem \ref{Basic25}, $I^{h}({\bf a})$ is a complete intersection.}
\end{ex1}

Finally, we provide a complete intersection criterion for $c_3<b_1+b_2$.

\begin{thm1} \label{Basic26} Suppose that $c_3<b_1+b_2$. Then $I^{h}({\bf a})$ is a complete intersection if and only if at least one of the following conditions holds: \begin{enumerate} \item[(A)] The element $c_{4}a_4 \in \mathbb{N}\{a_1,\ldots,a_4\}$ admits a factorization $c_{4}a_4=u_1a_{1}$ for some integer $u_{1}>c_1$, and the element $c_{3}a_3 \in \mathbb{N}\{a_1,\ldots,a_4\}$ admits a factorization $c_{3}a_3=v_1a_{1}+v_2a_{2}+v_4a_4$ for some $v_1, v_2, v_4 \in \mathbb{N}$ with $c_3 \geq v_1+v_2+v_4$.
\item[(B)] The element $c_{4}a_4 \in \mathbb{N}\{a_1,\ldots,a_4\}$ admits a factorization $c_{4}a_4=v_3a_{3}$ for some integer $v_{3}>c_3$, and the element $c_{3}a_3 \in \mathbb{N}\{a_1,\ldots,a_4\}$ admits a factorization $c_{3}a_3=u_1a_{1}$ for some $u_1 \in \mathbb{N}$ such that $u_1>c_3$ and $u_1>c_1$.
\end{enumerate}
\end{thm1}

{\em \noindent Proof.} $(\Leftarrow)$ Suppose first that condition (A) holds. Let $f_{1}=x_{2}^{c_2}-x_{1}^{c_1}, f_{2}=x_{1}^{u_1}-x_{4}^{c_4}$, and $f_{3}=x_{3}^{c_3}-x_{1}^{v_1}x_{2}^{v_2}x_4^{v_4}$, where $c_3 \geq v_1+v_2+v_4$. Consider the degree reverse lexicographic order $<$ with $x_1>x_2>x_3>x_4$. Then ${\rm in}_{<}(f_{1})=x_{2}^{c_2}$, ${\rm in}_{<}(f_{2})=x_{1}^{u_1}$, and ${\rm in}_{<}(f_{3})=x_{3}^{c_3}$ are pairwise relatively prime. Hence, by Remark \ref{VeryBasicRemark}(1), $\{f_{i} \mid 1 \leq i \leq 3\}$ is a Gr\"obner basis for $I({\bf a})$ with respect to $<$. By Remark \ref{VeryBasicRemark}(2), $\{f_{i}^h \mid 1 \leq i \leq 3\}$ is a Gr\"obner basis for $I^{h}({\bf a})$ with respect to $<_{0}$. Thus, $I^{h}({\bf a})$ is a complete intersection. 

Suppose now that condition (B) holds. Consider the binomials $f_{4}=x_{3}^{v_3}-x_{4}^{c_4} \in I({\bf a}), f_{5}=x_{1}^{u_1}-x_{3}^{c_3} \in I({\bf a})$, where $u_1 \geq c_3$. Then ${\rm in}_{<}(f_{1})$, ${\rm in}_{<}(f_{4})=x_{3}^{v_3}$, and ${\rm in}_{<}(f_{5})=x_{1}^{u_1}$ are pairwise relatively prime. Hence, by Remark \ref{VeryBasicRemark}(1), $\{f_1, f_4, f_5\}$ is a Gr\"obner basis for $I({\bf a})$ with respect to $<$. By Remark \ref{VeryBasicRemark}(2), $\{f_1^h, f_4^h, f_5^h\}$ is a Gr\"obner basis for $I^{h}({\bf a})$ with respect to $<_{0}$. Therefore, $I^{h}({\bf a})$ is a complete intersection.\\
$(\Rightarrow)$ Let $T=\{x_{1}^{u_1}-x_{0}^{u_0}x_{2}^{u_2}x_{3}^{u_3}x_{4}^{u_4}, x_{2}^{c_2}-x_{0}^{c_2-c_1}x_{1}^{c_1}, x_{3}^{v_3}-x_{0}^{v_0}x_{1}^{v_1}x_{2}^{v_2}x_{4}^{v_4}\}$ be a minimal generating set of $I^{h}({\bf a})$, where $u_{1} \geq u_2+u_3+u_4$ and $v_{3} \geq v_1+v_2+v_4$ with $v_3 \geq c_3$. Since $c_{4}<d_{1}+d_{2}+d_3$, the binomial $x_{0}^{d_{1}+d_{2}+d_3-c_{4}}x_{4}^{c_4}-x_{1}^{d_1}x_{2}^{d_2}x_{3}^{d_3}$ belongs to $I^{h}({\bf a})$. Since $T$ generates $I^{h}({\bf a})$, at least one of $x_{0}^{u_0}x_{2}^{u_2}x_{3}^{u_3}x_{4}^{u_4}$ and $x_{0}^{v_0}x_{1}^{v_1}x_{2}^{v_2}x_{4}^{v_4}$ divides $x_{0}^{d_{1}+d_{2}+d_3-c_{4}}x_{4}^{c_4}$. We distinguish two cases.\\ (1) Suppose that $x_{0}^{u_0}x_{2}^{u_2}x_{3}^{u_3}x_{4}^{u_4}$ divides $x_{0}^{d_{1}+d_{2}+d_3-c_{4}}x_{4}^{c_4}$. Then $u_{2}=u_{3}=0$ and $u_{4} \leq c_{4}$, so $x_{1}^{u_1}-x_{0}^{u_0}x_{4}^{u_4} \in I^{h}({\bf a})$, implying $x_{1}^{u_1}-x_{4}^{u_4} \in I({\bf a})$. Thus $u_{4} \geq c_{4}$, and therefore $u_{4}=c_{4}$. Since $c_1a_1 \neq c_4a_4$, we conclude that $u_1>c_1$. Additionally, $x_{0}^{b_1+b_2-c_3}x_{3}^{c_3}-x_{1}^{b_1}x_2^{b_2} \in I^{h}({\bf a})$. Since $T$ generates $I^{h}({\bf a})$, the monomial $x_3^{v_3}$ divides $x_{0}^{b_1+b_2-c_3}x_{3}^{c_3}$, and therefore $v_3 \leq c_3$. Thus $v_3=c_3$, so condition (A) holds.\\
(2) Suppose that $x_{0}^{v_0}x_{1}^{v_1}x_{2}^{v_2}x_{4}^{v_4}$ divides $x_{0}^{d_{1}+d_{2}+d_3-c_{4}}x_{4}^{c_4}$. Then $v_{1}=v_{2}=0$ and $v_{4} \leq c_{4}$, so $x_{3}^{v_3}-x_{0}^{v_0}x_{4}^{v_4} \in I^{h}({\bf a})$, implying $x_{3}^{v_3}-x_{4}^{v_4} \in I({\bf a})$. Thus $v_{4} \geq c_4$, and therefore $v_{4}=c_{4}$. Since $c_{3}a_{3} \neq c_4a_4$, it follows that $v_3>c_{3}$. Now, consider the binomial $x_{0}^{b_1+b_2-c_3}x_{3}^{c_3}-x_{1}^{b_1}x_2^{b_2} \in I^{h}({\bf a})$. Since $T$ generates $I^{h}({\bf a})$, the monomial $x_{0}^{u_0}x_{2}^{u_2}x_{3}^{u_3}x_{4}^{u_4}$ must divide $x_{0}^{b_1+b_2-c_3}x_{3}^{c_3}$, so $u_2=u_4=0$ and $u_3 \leq c_3$. Hence, $x_1^{u_1}-x_{0}^{u_0}x_{3}^{u_3} \in I^{h}({\bf a})$ with $u_1 \geq u_3$, implying $x_1^{u_1}-x_{3}^{u_3} \in I({\bf a})$ with $u_3 \geq c_3$. Thus $u_3=c_3$. Since $a_1 \neq a_3$, it follows that $u_1>c_3$, and because $c_1a_1 \neq c_3a_3$, we also have $u_1>c_1$. Thus condition (B) holds.  \hfill $\square$

\begin{rem1} {\rm Conditions {\rm (A)} and {\rm (B)} in Theorem~\ref{Basic26} can hold simultaneously. The sequence
${\bf a}=(a_1=3, a_2=12, a_3=8, a_4=48)$
provides such an example. The verification is analogous to that of Remark~\ref{rem:Basic24}.}
\end{rem1}

We conclude this section with an example illustrating condition {\rm (A)} in Theorem~\ref{Basic26}.

\begin{ex1} {\rm Let ${\bf a}=(a_{1}=3553,a_2=2431,a_3=3069,a_{4}=5168)$. Then $I({\bf a})$ is generated by $x_{2}^{19}-x_{1}^{13}$, $x_{3}^{17}-x_{1}x_{2}^{20}$, and $x_{4}^{11}-x_{1}^{3}x_{2}^{19}$. Here, $c_3=17<21=b_1+b_2$ and $c_4=11$. Since $c_4a_4$ does not admit a factorization $c_4a_4=v_3a_3$ with $v_3>c_3$, condition (B) of Theorem \ref{Basic26} fails. However, since $16a_1=11a_4$ and $17a_3=14a_1+a_2$, condition (A) of Theorem \ref{Basic26} is satisfied. Thus, $I^{h}({\bf a})$ is a complete intersection.}
\end{ex1}

\noindent \textbf{Conflict of interest statement}\\  The author declares no conflicts of interest.\\

\noindent \textbf{Data availability statement}\\Data sharing not applicable to this article as no datasets were generated or analysed during the current study.



\begin{thebibliography}{50}

\bibitem{BG2} I. Bermejo, I. Garc\'ia-Marco.
{\em Complete intersections in certain affine and projective monomial curves},
Bull. Braz. Math. Soc. (N.S.) {\bf 45} (2014), 599--624.

\bibitem{BG} I. Bermejo, I. Garc\'ia-Marco.
{\em Complete intersections in simplicial toric varieties},
J. Symbolic Comput. {\bf 68} (2015), 265--286.

\bibitem{Bresinsky75} H. Bresinsky.
{\em Symmetric semigroups of integers generated by 4 elements},
Manuscripta Math. {\bf 17} (3) (1975), 205--219.

\bibitem{CLO} D. Cox, J. Little, D. O’Shea.
{\em Ideals, Varieties, and Algorithms: An Introduction to Computational Algebraic Geometry and Commutative Algebra},
3rd ed., Springer, New York (2006).

\bibitem{Her} J. Herzog.
{\em Generators and relations of abelian semigroups and semigroup rings},
Manuscripta Math. {\bf 3} (1970), 175--193.

\bibitem{HerSta} J. Herzog, D. Stamate.
{\em Cohen--Macaulay criteria for projective monomial curves via Gr\"obner bases},
Acta Math. Vietnam. {\bf 44} (2019), 51--64.

\bibitem{Katsa} A. Katsabekis.
{\em Complete intersection monomial curves and the Cohen--Macaulayness of their tangent cones},
Algebra Colloq. {\bf 26} (2019), no. 4, 629--642.

\bibitem{Katsa1} A. Katsabekis.
{\em Projective closure of Gorenstein monomial curves and the Cohen--Macaulay property},
Arch. Math. {\bf 122} (2024), no. 3, 273--283.

\bibitem{KO} A. Katsabekis, I. Ojeda.
{\em An indispensable classification of monomial curves in $\mathbb{A}^{4}(\mathbb{K})$},
Pacific J. Math. {\bf 268} (2014), 95--116.

\bibitem{PatilRoberts2003} D. P. Patil, L. G. Roberts. {\em Hilbert functions of monomial curves}, J. Pure Appl. Algebra {\bf 183} (2003), 275--292.

\bibitem{SSS} J. Saha, I. Sengupta, P. Srivastava.
{\em Betti sequence of the projective closure of affine monomial curves},
J. Symbolic Comput. {\bf 119} (2023), 101--111.

\bibitem{Shi} T. Shibuta.
{\em Cohen--Macaulayness of almost complete intersection tangent cones},
J. Algebra {\bf 319} (2008), 3222--3243.

\bibitem{Sturmfels95} B. Sturmfels.
{\em Gr\"obner Bases and Convex Polytopes},
University Lecture Series, Vol. 8,
American Mathematical Society, Providence, RI, 1996.

	
\end{thebibliography}


\end{document}

\section{Projective closures in case 1}

In this section, we study in detail the first case arising from the structure of the minimal generators. For any $f \in K[x_1,\ldots,x_n]$, we denote by $f^{h}$ its homogenization with respect to $x_0$. The following remark will be used repeatedly.

\begin{rem1} \label{VeryBasicRemark} 
{\rm Let $I \subset K[x_1, \ldots, x_n]$ be an ideal, and let $<$ be a graded monomial order on $K[x_1, \ldots, x_n]$. If $\{f_1, \ldots, f_r\}$ is a Gr\"obner basis of $I$ with respect to $<$, then the set $\{f_1^h, \ldots, f_r^h\}$ is a Gr\"obner basis of $I^h \subset K[x_0, x_1, \ldots, x_n]$ with respect to the graded monomial order $<_0$ that extends $<$ and satisfies $x_0 <_{0} x_i$ for all $i \geq 1$. That is, $<_0$ preserves the relative order of $x_1, \ldots, x_n$ and places $x_0$ below all other variables (see \cite[Chapter 8]{CLO}).
}
\end{rem1}

Let ${\bf a}=(a_1,\ldots,a_4)$ be a sequence of distinct relatively prime positive integers. 

\begin{not1} 
{\rm Throughout the paper, let $c_i$ denote the smallest positive integer such that $c_{i} a_{i} \in \sum_{j \neq i} \mathbb{N} a_{j}$, for each $i=1,\ldots,4$, where $\mathbb{N}$ is the set of nonnegative integers.}
\end{not1}

\begin{thm1} \label{CasesMinimal} 
Let $I({\bf a})$ be a complete intersection toric ideal, where ${\bf a}=(a_1,\ldots,a_4)$ is a sequence of distinct positive integers with ${\rm gcd}(a_1,\ldots,a_4)=1$. Then after permuting the variables, if necessary, there exists a minimal system of binomial generators $S$ of $I({\bf a})$ of the following form:
\begin{enumerate}
    \item[{\rm CASE 1:}] If $c_{1}a_{1}=c_{2}a_{2}$ and $c_{3}a_{3}=c_{4}a_{4}$, then 
    $$S=\{x_{2}^{c_2}-x_{1}^{c_1}, x_{3}^{c_3}-x_{4}^{c_4}, x_{1}^{d_1}x_{2}^{d_2}-x_{3}^{d_3}x_{4}^{d_4}\}.$$
    
    \item[{\rm CASE 2:}] If $c_{1}a_{1}=c_{2}a_{2}=c_{3}a_{3} \neq c_{4}a_{4}$, then 
    $$S=\{x_{2}^{c_2}-x_{1}^{c_1}, x_{3}^{c_3}-x_{1}^{c_1}, x_{4}^{c_4}-x_{1}^{d_1}x_{2}^{d_2}x_{3}^{d_3}\},$$ 
    where at least two of $d_1$, $d_2$, and $d_3$ are nonzero.
    
    \item[{\rm CASE 3:}] If $c_{1}a_{1}=c_{2}a_{2}$ and $c_{i}a_{i} \neq c_{j}a_{j}$ for all $\{i,j\} \neq \{1,2\}$, then 
    $$S=\{x_{2}^{c_2}-x_{1}^{c_1}, x_{3}^{c_3}-x_{1}^{b_1}x_{2}^{b_2}, x_{4}^{c_4}-x_{1}^{d_1}x_{2}^{d_2}x_{3}^{d_3}\},$$ 
    where both $b_1$ and $b_2$ are nonzero, and at least two of $d_1$, $d_2$, and $d_3$ are nonzero.
\end{enumerate}
\end{thm1}

{\em \noindent Proof.} We begin by noting that if $I({\bf a})$ is a complete intersection, then one of the following holds: either $c_1a_1=c_2a_2$ and $c_3a_3=c_4a_4$; or $c_{1}a_{1}=c_{2}a_{2}=c_{3}a_{3} \neq c_{4}a_{4}$; or $c_{1}a_{1}=c_{2}a_{2}$ and $c_{i}a_{i} \neq c_{j}a_{j}$ for all $\{i,j\} \neq \{1,2\}$ (see \cite{KO}). By \cite[Proposition 3.2]{Shi}, a minimal generating set of $I({\bf a})$ is of the form 
$$\{x_{2}^{c_2}-x_{1}^{c_1}, x_{3}^{c_3}-x_{4}^{c_4}, x_{1}^{d_1}x_{2}^{d_2}-x_{3}^{d_3}x_{4}^{d_4}\}$$ 
in Case 1, and 
$$\{x_{2}^{c_2}-x_{1}^{c_1}, x_{3}^{c_3}-x_{1}^{b_1}x_{2}^{b_2}, x_{4}^{c_4}-x_{1}^{d_1}x_{2}^{d_2}x_{3}^{d_3}\}$$ 
in Cases 2 and 3. In Case 2, since $c_{1}a_1=c_{2}a_2=c_{3}a_{3}$, we can set $b_{1}=c_1$ and $b_2=0$, so that the second generator is $x_{3}^{c_3}-x_{1}^{c_1}$. We may also assume that at least two of $d_1$, $d_2$, and $d_3$ are nonzero; otherwise, if, for instance, $d_{1}=d_{2}=0$, then $d_{3}>c_3$, and using $x_{3}^{c_3}-x_{1}^{c_1}$, we can replace $x_{4}^{c_4}-x_{3}^{d_3}$ with $x_{4}^{c_4}-x_{1}^{c_1}x_{3}^{d_3-c_3} \in I({\bf a})$. In Case 3, we may assume that both $b_1$ and $b_2$ are nonzero; indeed, if, for instance, $b_2=0$, then $b_1>c_1$, and using $x_2^{c_2}-x_1^{c_1}$ we can replace $x_3^{c_3} - x_{1}^{b_1}x_{2}^{b_2}$ with $x_3^{c_3}-x_{1}^{b_1-c_1}x_{2}^{b_2+c_2} \in I({\bf a})$. Similarly, we may assume that at least two of $d_1$, $d_2$, and $d_3$ are nonzero. \hfill $\square$ 

For the rest of this section, let $I({\bf a})$ be a complete intersection toric ideal minimally generated by $x_{2}^{c_2}-x_{1}^{c_1}$, $x_{3}^{c_3}-x_{4}^{c_4}$, and $x_{1}^{d_1}x_{2}^{d_2}-x_{3}^{d_3}x_{4}^{d_4}$. Without loss of generality, we fix the assumptions $a_{4}={\rm max}\{a_1,\ldots,a_4\}$, $c_2>c_1$, $d_2 \leq c_2$, and $d_3 \leq c_3$. Indeed, if $d_2>c_2$, we may write $d_{2}=kc_2+l$ for some integers $k \geq 1$ and $0 \leq l<c_2$, and replace $x_{1}^{d_1}x_{2}^{d_2}-x_{3}^{d_3}x_{4}^{d_4}$ with $x_{1}^{kc_1+d_1}x_{2}^{l}-x_{3}^{d_3}x_{4}^{d_4} \in I({\bf a})$. Similarly, we may assume $d_3 \leq c_3$. Finally, since $a_4>a_3$, we have $c_4<c_3$.

Throughout this section, let $<$ be the degree reverse lexicographic order on $K[x_1, \ldots, x_4]$ with $x_1 > x_2 > x_3 > x_4$. We begin by considering a first subcase in which the ideal $I^{h}({\bf a})$ is a complete intersection.

\begin{thm1} \label{BAS1} If $d_2=c_2$ and $d_3=c_3$, then $I^{h}({\bf a})$ is a complete intersection.
\end{thm1}

{\em \noindent Proof.} Let $f_1=x_2^{c_2}-x_1^{c_1}$, $f_2=x_3^{c_3}-x_4^{c_4}$, and $g = x_1^{c_1+d_1}-x_4^{c_4+d_4} \in I({\bf a})$ since $d_2=c_2$ and $d_3=c_3$. Their initial monomials $x_2^{c_2}$, $x_3^{c_3}$, and $x_1^{c_1 + d_1}$ are pairwise relatively prime, and hence $\{f_1, f_2, g\}$ is a Gr\"obner basis for $I({\bf a})$ with respect to $<$. By Remark \ref{VeryBasicRemark}, the set $\{f_1^{h}, f_2^{h}, g^{h} \}$ is a Gr\"obner basis for $I^h({\bf a})$ with respect to $<_{0}$. Hence, $I^h({\bf a})$ is a complete intersection.

\hfill $\square$ \par\vspace{0.5\baselineskip}

The following example illustrates the subcase considered in Theorem~\ref{BAS1}.

\begin{ex1} {\rm Let ${\bf a}=(132, 60, 165, 220)$. Then $I({\bf a})$ is generated by $x_{2}^{11}-x_{1}^{5}$, $x_3^{4}-x_{4}^{3}$, and $x_{2}^{11}-x_{3}^{4}$. Here, $d_{2}=c_2=11$ and $d_{3}=c_3=4$. By Theorem \ref{BAS1}, the ideal $I^{h}({\bf a})$ is a complete intersection.}
\end{ex1}

Next, we consider another subcase in which $I^{h}({\bf a})$ is a complete intersection.

\begin{thm1} \label{CI23} Suppose that $d_2=0$ and $d_{3}=c_{3}$. Then $I^{h}({\bf a})$ is a complete intersection.
\end{thm1}

{\em \noindent Proof.} Let $f_1 = x_2^{c_2} - x_1^{c_1}$, $f_2 = x_3^{c_3} - x_4^{c_4}$, and $g = x_1^{d_1} - x_4^{c_4 + d_4} \in I({\bf a})$. The initial monomials $x_2^{c_2}$, $x_3^{c_3}$, and $x_1^{d_1}$ are pairwise relatively prime; hence, $\{f_1, f_2, g\}$ is a Gr\"obner basis for $I({\bf a})$ with respect to $<$. By Remark \ref{VeryBasicRemark}, $\{f_1^{h}, f_2^{h}, g^{h} \}$ is a Gr\"obner basis for $I^h({\bf a})$ with respect to $<_{0}$. Thus, $I^h({\bf a})$ is a complete intersection.

\hfill $\square$ \par\vspace{0.5\baselineskip}

We now obtain several sufficient and necessary conditions under which the arithmetically Cohen--Macaulay property holds.

\begin{thm1} \label{ACM1} Suppose that $d_2=c_2$ and $d_3<c_3$. Then $\overline{C({\bf a})}$ is arithmetically Cohen--Macaulay if and only if $c_{1}+d_{1} \geq d_{3}+d_{4}$.
\end{thm1}

{\em \noindent Proof.} ($\Leftarrow$) Consider the set $T=\{f_{1}=x_2^{c_2}-x_1^{c_1}, f_{2}=x_3^{c_3}-x_4^{c_4}, f_{3}=x_1^{c_1+d_{1}}-x_{3}^{d_3}x_4^{d_4}\} \subset I({\bf a})$. Then ${\rm in}_{<}(f_{1})=x_2^{c_2}$, ${\rm in}_{<}(f_{2})=x_3^{c_3}$, and ${\rm in}_{<}(f_{3})=x_1^{c_1 + d_1}$ are pairwise relatively prime. Thus $T$ is a Gr\"obner basis for $I({\bf a})$ with respect to $<$. Since $x_4$ does not divide the initial monomial of any element in $T$, it follows from \cite[Theorem 2.2]{HerSta} that $\overline{C({\bf a})}$ is arithmetically Cohen--Macaulay.

\hfill $\square$ \par\vspace{0.5\baselineskip}

As an application of Theorem \ref{ACM1}, we obtain the following corollary. 

\begin{cor1} \label{CorCase1CI} Suppose that $d_2=c_2$ and $d_{3}<c_{3}$. Then the ideal $I^{h}({\bf a})$ is a complete intersection if and only if $c_{1}+d_{1} \geq d_{3}+d_{4}$.
\end{cor1}

{\em \noindent Proof.} ($\Leftarrow$) By Remark \ref{VeryBasicRemark}, the set $\{f_{i}^{h} \mid 1 \leq i \leq 3\}$ is a Gr\"obner basis for $I^{h}({\bf a})$ with respect to $<_{0}$; hence $I^{h}({\bf a})$ is a complete intersection.\\

($\Rightarrow$) If $I^{h}({\bf a})$ is a complete intersection, then $\overline{C({\bf a})}$ is arithmetically Cohen--Macaulay. By Theorem \ref{ACM1}, it follows that $c_1+d_{1} \geq d_{3}+d_{4}$.

\hfill $\square$ \par\vspace{0.5\baselineskip}

\begin{thm1} \label{ACM2} Suppose that $0<d_2<c_2$ and $d_{3}=c_{3}$. Then $\overline{C({\bf a})}$ is arithmetically Cohen--Macaulay if and only if $c_{1}+d_{1} \geq c_{2}-d_{2}+c_{4}+d_{4}$.
\end{thm1}

{\em \noindent Proof.} ($\Leftarrow$) Consider the set $T=\{f_1=x_2^{c_2}-x_1^{c_1},\; f_2=x_3^{c_3}-x_4^{c_4},\; f_3=x_1^{d_1}x_2^{d_2}-x_4^{c_4+d_4},\; f_4=x_1^{c_1+d_1}-x_2^{c_{2}-d_2}x_4^{c_4+d_4}\} \subset I({\bf a})$. Since $c_2>c_1$ and $c_1+d_1 \geq c_2-d_2+c_4+d_4$, the initial monomials with respect to $<$ are ${\rm in}_{<}(f_{1})=x_2^{c_2}$, ${\rm in}_{<}(f_{2})=x_3^{c_3}$, ${\rm in}_{<}(f_{3})=x_1^{d_1}x_2^{d_2}$, and ${\rm in}_{<}(f_4)=x_1^{c_1+d_{1}}$. Since ${\rm in}_{<}(f_1)$ and ${\rm in}_{<}(f_2)$ are relatively prime, it follows that $S(f_{1},f_{2}) \stackrel{T}{\longrightarrow} 0$. Similarly, $S(f_{1},f_4) \stackrel{T}{\longrightarrow} 0$, $S(f_{2},f_{3}) \stackrel{T}{\longrightarrow} 0$, and $S(f_{2},f_4) \stackrel{T}{\longrightarrow} 0$. We have that $S(f_{1},f_{3})=-f_4 \stackrel{f_4}{\longrightarrow} 0$ and $S(f_{3},f_4)=x_4^{c_4+d_4}(x_{2}^{c_2}-x_{1}^{c_1}) \stackrel{f_1}{\longrightarrow} 0$. Thus $T$ is a Gr\"obner basis for $I({\bf a})$ with respect to $<$. Since none of the initial monomials is divisible by $x_4$, it follows from \cite[Theorem 2.2]{HerSta} that $\overline{C(\mathbf a)}$ is arithmetically Cohen--Macaulay.\\
($\Rightarrow$) Assume that $c_{1}+d_{1}<c_{2}-d_{2}+c_{4}+d_{4}$. Since $\overline{C({\bf a})}$ is arithmetically Cohen--Macaulay and ${\rm in}_{<}(f_4)=x_2^{c_2-d_2}x_4^{c_4+d_4}$, by \cite[Theorem 2.2]{HerSta} there exists a binomial $B \in I({\bf a})$ such that ${\rm in}_{<}(B)$ divides $x_2^{c_{2}-d_2}$, contradicting the minimality of $c_2$. \hfill $\square$

We conclude this section with an example illustrating the complete intersection criterion given in Corollary~\ref{CorCase1CI}.

\begin{ex1} {\rm Let ${\bf a}=(14, 7, 8, 20)$ and ${\bf v}=(7,7,4,10)$. By \cite[Theorem 2.6]{Katsa}, for any $t \in \mathbb{N}$ such that ${\rm gcd}(14+7t,7+7t,8+4t,20+10t)=1$, the ideal $I({\bf a}+t{\bf v})$ is generated by $x_{2}^{2+t}-x_{1}^{1+t}$, $x_{3}^{5}-x_{4}^{2}$, and $x_{1}^{2}-x_{3}x_{4}$. Here, $d_1=2$, $d_2=0$, $d_3=d_4=1$, and $c_3=5$. By Corollary \ref{CorCase1CI}, the ideal $I^{h}({\bf a}+t{\bf v})$ is a complete intersection whenever ${\rm gcd}(14+7t,7+7t,8+4t,20+10t)=1$.}
\end{ex1}

\section{Projective closures in case~2}

This section focuses on Case 2 of Theorem \ref{CasesMinimal}, where $I({\bf a})$ has a minimal generating set $\{g_{1}=x_{2}^{c_2}-x_{1}^{c_1}, g_{2}=x_{3}^{c_3}-x_{1}^{c_1}, g_{3}=x_{4}^{c_4}-x_{1}^{d_1}x_{2}^{d_2}x_{3}^{d_3}\}$, and at least two of $d_1$, $d_2$, and $d_3$ are nonzero.

We begin by studying the case where $a_{i} = {\rm max}\{a_1,\ldots,a_4\}$ for some $1 \leq i \leq 3$. Without loss of generality, assume $a_1={\rm max}\{a_1,\ldots,a_4\}$, and fix $d_{2} \leq c_2$ and $d_{3} \leq c_3$.

Throughout this section, unless stated otherwise, let $<$ denote the degree reverse lexicographic order on $K[x_1,\ldots,x_4]$ with $x_4>x_3>x_2>x_1$.

We first consider the case $d_2=c_{2}$ and $d_{3}=c_3$.

\begin{thm1} \label{Veryfirst} 
If $d_2=c_{2}$ and $d_{3}=c_3$, then $I^{h}({\bf a})$ is a complete intersection.
\end{thm1}

{\em \noindent Proof.} Using $g_1$ and $g_2$, we obtain $f=x_{4}^{c_4}-x_{1}^{d_1+2c_1} \in I({\bf a})$. Then ${\rm in}_{<}(g_1)=x_{2}^{c_2}$, ${\rm in}_{<}(g_2)=x_{3}^{c_3}$, and ${\rm in}_{<}(f)=x_{4}^{c_4}$ are pairwise relatively prime. Hence, $\{g_1, g_2, f\}$ is a Gr\"obner basis for $I({\bf a})$ with respect to $<$. By Remark \ref{VeryBasicRemark}, $\{g_{1}^{h}, g_{2}^{h}, f^{h}\}$ is a Gr\"obner basis for $I^{h}({\bf a})$ with respect to $<_{0}$, so $I^{h}({\bf a})$ is a complete intersection. \hfill $\square$ \par\vspace{0.5\baselineskip}

To illustrate Theorem \ref{Veryfirst}, we present the following example.

\begin{ex1} 
{\rm Let ${\bf a}=(6851, 5797, 4433, 4862)$. Then $I({\bf a})$ is generated by $x_{2}^{13}-x_1^{11}$, $x_{3}^{17}-x_{1}^{11}$, and $x_{4}^{31}-x_{2}^{13}x_{3}^{17}$. Here, $c_2=d_2=13$ and $c_3=d_3=17$. By Theorem \ref{Veryfirst}, $I^{h}({\bf a})$ is a complete intersection.}
\end{ex1}

Next, we examine the case $d_2<c_2$ and $d_3=c_3$.

\begin{thm1} \label{ACM5} 
If $d_2<c_{2}$ and $d_{3}=c_3$, then $\overline{C({\bf a})}$ is arithmetically Cohen--Macaulay if and only if $c_{4} \geq c_1+d_{1}+d_{2}$.
\end{thm1}

{\em \noindent Proof.} 
($\Leftarrow$) Let $f=x_{4}^{c_4}-x_{1}^{d_1+c_1}x_{2}^{d_2} \in I({\bf a})$. Since $c_{4} \geq c_1+d_{1}+d_{2}$, we have ${\rm in}_{<}(f)=x_{4}^{c_4}$. Then ${\rm in}_{<}(g_1)$, ${\rm in}_{<}(g_2)$, and ${\rm in}_{<}(f)$ are pairwise relatively prime, so $\{g_1, g_2, f\}$ is a Gr\"obner basis for $I({\bf a})$ with respect to $<$. As $x_1$ divides none of these initial monomials, $\overline{C({\bf a})}$ is arithmetically Cohen--Macaulay by \cite[Theorem 2.2]{HerSta}.\\
($\Rightarrow$) Assume that $c_{4}<c_1+d_{1}+d_{2}$. Since $\overline{C({\bf a})}$ is arithmetically Cohen--Macaulay and ${\rm in}_{<}(f)=x_{1}^{d_1+c_1}x_{2}^{d_2}$, \cite[Theorem~2.2]{HerSta} guarantees the existence of a binomial $B \in I({\bf a})$ such that ${\rm in}_{<}(B)$ divides $x_{2}^{d_2}$. Thus ${\rm in}_{<}(B)=x_{2}^{e}$ for some $e \leq d_2$, which implies $e<c_2$, contradicting the minimality of $c_2$. \hfill $\square$ \par\vspace{0.5\baselineskip}

The following corollary is a consequence of Theorem \ref{ACM5}.

\begin{cor1} \label{Cor12} 
If $d_2<c_{2}$ and $d_{3}=c_3$, then $I^{h}({\bf a})$ is a complete intersection if and only if $c_{4} \geq c_1+d_{1}+d_{2}$.
\end{cor1}

{\em \noindent Proof.} If $I^{h}({\bf a})$ is a complete intersection, then $\overline{C({\bf a})}$ is arithmetically Cohen--Macaulay, and hence $c_{4} \geq c_1+d_{1}+d_{2}$ by Theorem \ref{ACM5}. Conversely, assume that $c_{4} \geq c_1+d_{1}+d_{2}$. By Remark \ref{VeryBasicRemark}, the set $\{g_{1}^{h}, g_{2}^{h}, f^{h}\}$ is a Gr\"obner basis for $I^{h}({\bf a})$ with respect to $<_{0}$. Consequently, $I^{h}({\bf a})$ is a complete intersection. \hfill $\square$ \par\vspace{0.5\baselineskip}

The following example illustrates Corollary \ref{Cor12}.

\begin{ex1} 
{\rm Let ${\bf a}=(60, 20, 15, 16)$. Then $I({\bf a})$ is generated by $x_{2}^3-x_1$, $x_{3}^{4}-x_{1}$, and $x_{4}^{5}-x_{2}x_{3}^4$. Here, $d_2=1<3=c_2$ and $d_3=c_3=4$. Since $c_4=5>2=c_1+d_1+d_2$, Corollary~\ref{Cor12} implies that $I^{h}({\bf a})$ is a complete intersection.}
\end{ex1}

We now consider the case $d_2=c_2$ and $d_3<c_3$.

\begin{thm1} \label{ACM6} 
If $d_2=c_{2}$ and $d_{3}<c_3$, then $\overline{C({\bf a})}$ is arithmetically Cohen--Macaulay if and only if $c_{4} \geq d_1+c_1+d_{3}$.
\end{thm1}

{\em \noindent Proof.} 
($\Leftarrow$) Let $f=x_{4}^{c_4}-x_{1}^{d_1+c_1}x_{3}^{d_3} \in I(\mathbf a)$. Since $c_{4} \geq d_1+c_1+d_{3}$, we have ${\rm in}_{<}(f)=x_{4}^{c_4}$. Then ${\rm in}_{<}(g_1)$, ${\rm in}_{<}(g_2)$, and ${\rm in}_{<}(f)$ are pairwise relatively prime, so $\{g_1, g_2, f\}$ is a Gr\"obner basis for $I({\bf a})$ with respect to $<$. As $x_1$ divides none of these initial monomials, $\overline{C({\bf a})}$ is arithmetically Cohen--Macaulay by \cite[Theorem 2.2]{HerSta}.\\
($\Rightarrow$) Assume that $c_{4}<c_1+d_{1}+d_{3}$. Since $\overline{C({\bf a})}$ is arithmetically Cohen--Macaulay and ${\rm in}_{<}(f)=x_{1}^{d_1+c_1}x_{3}^{d_3}$, \cite[Theorem 2.2]{HerSta} guarantees the existence of a binomial $B \in I({\bf a})$ such that ${\rm in}_{<}(B)$ divides $x_{3}^{d_3}$. Thus ${\rm in}_{<}(B)=x_{3}^{e}$ for some $e \leq d_3$, which implies $e<c_3$, contradicting the minimality of $c_3$. \hfill $\square$ \par\vspace{0.5\baselineskip}

Theorem \ref{ACM6} yields the following complete intersection criterion.

\begin{cor1} \label{Cor13} 
If $d_2=c_{2}$ and $d_{3}<c_3$, then $I^{h}({\bf a})$ is a complete intersection if and only if $c_{4} \geq c_1+d_{1}+d_{3}$.
\end{cor1}

{\em \noindent Proof.} If $I^{h}({\bf a})$ is a complete intersection, then $\overline{C({\bf a})}$ is arithmetically Cohen--Macaulay. By Theorem~\ref{ACM6}, it follows that $c_4 \ge c_1+d_1+d_3$. Conversely, assume that $c_{4} \geq c_1+d_{1}+d_{3}$. By Remark \ref{VeryBasicRemark}, the set $\{g_{1}^{h}, g_{2}^{h}, f^{h}\}$ is a Gr\"obner basis of $I^{h}({\bf a})$ with respect to $<_{0}$. Consequently, $I^{h}({\bf a})$ is a complete intersection. \hfill $\square$ \par\vspace{0.5\baselineskip}

We now study the case $d_2<c_2$ and $d_3<c_3$.

\begin{thm1} \label{Basic4} 
Suppose that $d_2<c_2$ and $d_{3}<c_3$. Then the ideal $I^{h}({\bf a})$ is a complete intersection if and only if $c_{4} \geq d_{1}+d_{2}+d_{3}$.
\end{thm1}

{\em \noindent Proof.} ($\Leftarrow$) Since $c_{4} \geq d_1+d_2+d_{3}$, we have ${\rm in}_{<}(g_3)=x_{4}^{c_4}$. Then ${\rm in}_{<}(g_1)$, ${\rm in}_{<}(g_2)$, and ${\rm in}_{<}(g_3)$ are pairwise relatively prime, so $\{g_{i} \mid 1 \leq i \leq 3\}$ is a Gr\"obner basis of $I({\bf a})$ with respect to $<$. By Remark \ref{VeryBasicRemark}, $\{g_{i}^{h} \mid 1 \leq i \leq 3\}$ is a Gr\"obner basis for $I^{h}({\bf a})$ with respect to $<_{0}$. Therefore, $I^{h}({\bf a})$ is a complete intersection.\\ 
($\Rightarrow$) Let $T=\{g_{1}^{h}, g_{2}^{h}, x_{4}^{u_4}-x_{0}^{u_0}x_{1}^{u_1}x_{2}^{u_2}x_{3}^{u_3}\}$ be a minimal generating set of $I^{h}({\bf a})$, where $u_{4} \geq u_{1}+u_{2}+u_{3}$. Note that $u_{4} \geq c_{4}$ since $x_{4}^{u_4}-x_{1}^{u_1}x_{2}^{u_2}x_{3}^{u_3} \in I({\bf a})$. Assume that $c_{4}<d_1+d_{2}+d_{3}$. Since $g_{3}^{h}=x_{0}^{d_{1}+d_{2}+d_{3}-c_{4}}x_{4}^{c_{4}}-x_{1}^{d_{1}}x_{2}^{d_{2}}x_{3}^{d_{3}} \in I^{h}({\bf a})$ and the set $T$ generates $I^{h}({\bf a})$, the monomial $x_{4}^{u_4}$ divides $x_{4}^{c_4}$, so $u_4 \leq c_4$, and hence $u_4=c_4$. Moreover, since $d_2<c_2$ and $d_3<c_3$, $x_{1}^{u_1}x_{2}^{u_2}x_{3}^{u_3}$ divides $x_{1}^{d_1}x_{2}^{d_{2}}x_{3}^{d_{3}}$, so $u_i \leq d_i$ for each $1 \leq i \leq 3$. Consequently, $x_{4}^{c_4} (x_1^{d_1-u_1}x_2^{d_2-u_2}x_3^{d_3-u_3}-1) \in I({\bf a})$, which implies that $x_1^{d_1-u_1}x_2^{d_2-u_2}x_3^{d_3-u_3}-1 \in I({\bf a})$, since $I({\bf a})$ is prime and does not contain the variable $x_4$. Therefore, $d_i=u_i$ for all $1 \leq i \leq 3$, and hence $c_4 \geq d_1+d_2+d_3$, a contradiction. \hfill $\square$ \par\vspace{0.5\baselineskip}

When $d_{2}<c_2$, $d_3<c_3$, and $c_4<d_1+d_2+d_3$, the ideal $I^{h}({\bf a})$ is not a complete intersection. Nevertheless, the following Gr\"obner basis computations show that $\overline{C({\bf a})}$ may still be arithmetically Cohen--Macaulay.

Let $w$ be the smallest positive integer such that $c_2-wd_{2} \leq 0$ or $c_3-wd_{3} \leq 0$, namely $w={\rm min} \{l \in \mathbb{N} \setminus \{0\} \mid c_{2}-ld_{2} \leq 0 \ \textrm{or} \ c_{3}-ld_{3} \leq 0\}.$ Since $d_2<c_2$ and $d_3<c_3$, we have $w \geq 2$. For $0 \leq i \leq w-2$ define 
$$p_{i}=x_{2}^{c_{2}-(i+1)d_{2}}x_{4}^{(i+1)c_{4}}-x_{1}^{c_1}x_{3}^{(i+1)d_{3}}, \ q_{i}=x_{3}^{c_{3}-(i+1)d_3}x_{4}^{(i+1)c_{4}}-x_{1}^{c_1}x_{2}^{(i+1)d_2}.$$ 
Additionally, define 
\[
r =
\begin{cases}
x_4^{wc_4}-x_1^{c_1}x_2^{wd_2}x_3^{wd_3-c_3},
& \text{if } c_2-wd_2>0,\\[4pt]
x_4^{wc_4}-x_1^{c_1}x_2^{wd_2-c_2}x_3^{wd_3},
& \text{if } c_2-wd_2 \leq 0.
\end{cases}
\]

\begin{prop1} \label{CMBasic12} 
Let $d_1=0$, $d_{2}<c_{2}$, $d_{3}<c_{3}$, $c_3>c_2$, and $c_4<d_{2}+d_{3}$. Suppose that the following conditions hold:
\begin{enumerate} 
\item $(w-1)(c_{4}-d_{2}-d_{3})+c_{2}-c_{1} \geq 0$.
\item $w(c_{4}-d_{2}-d_{3})+c_{3}-c_{1} \geq 0$ when $c_{2}-wd_{2}> 0$.
\item $w(c_{4}-d_{2}-d_{3})+c_{2}-c_{1} \geq 0$ when $c_{2}-wd_{2} \leq 0.
\end{enumerate}
Then the set 
\[
\begin{aligned}
T = {} & \{\, g_{1}=x_{2}^{c_2}-x_{1}^{c_1},\;
          g_{2}=x_{3}^{c_3}-x_{1}^{c_1},\;
          f=x_{2}^{d_2}x_{3}^{d_3}-x_{4}^{c_4} \,\} \cup \{\, p_i \mid 0 \le i \le w-2 \,\} \\
      & 
        \cup \{\, q_i \mid 0 \le i \le w-2 \,\}
        \cup \{\, r \,\}.
\end{aligned}
\]
is a Gr\"obner basis for $I({\bf a})$ with respect to $<$.
\end{prop1}

{\em \noindent Proof.} For each $0 \leq i \leq w-2$, the binomials $p_{i}$, $q_{i}$, and $r$ belong to the ideal $I({\bf a})$ by direct computation. Here, ${\rm in}_{<}(g_{1})=x_{2}^{c_{2}}$, ${\rm in}_{<}(g_{2})=x_{3}^{c_{3}}$, and ${\rm in}_{<}(f)=x_{2}^{d_{2}}x_{3}^{d_{3}}$. Since $i+1 \leq w-1$ for all $0 \leq i \leq w-2$, it follows that $(i+1)(c_{4}-d_{2}-d_{3})+c_{2}-c_{1} \geq (w-1)(c_{4}-d_{2}-d_{3})+c_{2}-c_{1}$, as $c_{4}-d_{2}-d_{3}<0$. Hence, $(i+1)(c_{4}-d_{2}-d_{3})+c_{2}-c_{1} \geq 0$, which implies that ${\rm in}_{<}(p_{i})=x_{2}^{c_{2}-(i+1)d_{2}}x_{4}^{(i+1)c_{4}}$. Since $c_3>c_2$, it follows that $(w-1)(c_{4}-d_{2}-d_{3})+c_{3}-c_{1}>(w-1)(c_{4}-d_{2}-d_{3})+c_{2}-c_{1}$. Thus $(w-1)(c_{4}-d_{2}-d_{3})+c_{3}-c_{1} >0$, and therefore ${\rm in}_{<}(q_{i})=x_{3}^{c_{3}-(i+1)d_3}x_{4}^{(i+1)c_{4}}$. Finally, ${\rm in}_{<}(r)=x_{4}^{wc_4}$. Since ${\rm in}_{<}(g_{1})$ and ${\rm in}_{<}(g_{2})$ are relatively prime, we deduce that $S(g_{1},g_{2}) \stackrel{T} {\longrightarrow} 0$. Similarly, we have $S(g_{1},q_{i}) \stackrel{T} {\longrightarrow} 0$ and $S(g_{2},p_i) \stackrel{T} {\longrightarrow} 0$ for all $0 \leq i \leq w-2$, as well as $S(g_{1},r) \stackrel{T} {\longrightarrow} 0$, $S(g_{2},r) \stackrel{T} {\longrightarrow} 0$, and $S(f,r) \stackrel{T} {\longrightarrow} 0$. 

We have that $S(g_{1},f)=p_{0} \stackrel{p_{0}} {\longrightarrow} 0$ and $S(g_{2},f)=q_{0} \stackrel{q_{0}} {\longrightarrow} 0$. We now show that $S(g_{1}, p_{i}) \stackrel{T}{\longrightarrow} 0$ for every $0 \leq i \leq w-2$. Observe that $$S(g_{1},p_{i})=x_{1}^{c_1}(x_{2}^{(i+1)d_2}x_{3}^{(i+1)d_3}-x_{4}^{(i+1)c_4}) \stackrel{f} {\longrightarrow} 0.$$ Similarly, we have $S(g_{2},q_{i})=x_{1}^{c_1}(x_{2}^{(i+1)d_2}x_{3}^{(i+1)d_3}-x_{4}^{(i+1)c_4}) \stackrel{f} {\longrightarrow} 0.$ 

For each $0 \leq i \leq w-3$ we have $S(f,p_{i})=x_{1}^{c_1}x_{3}^{(i+2)d_3}-x_{2}^{c_2-(i+2)d_2}x_{4}^{(i+2)c_{4}} \stackrel{p_{i+1}} {\longrightarrow} 0$. If $c_{2}-wd_{2} \leq 0$, then $S(f,p_{w-2})=x_{1}^{c_1}x_{2}^{wd_{2}-c_{2}}x_{3}^{wd_3}-x_{4}^{wc_4} \stackrel{r} {\longrightarrow} 0$. On the other hand, if $c_{2}-wd_{2}>0$, then $S(f,p_{w-2})=x_{1}^{c_1}x_{3}^{wd_3}-x_{2}^{c_{2}-wd_2}x_4^{wc_4} \stackrel{r} {\longrightarrow} x_{1}^{c_1}x_{3}^{wd_3-c_3}(x_{3}^{c_3}-x_{2}^{c_{2}}) \stackrel{g_1} {\longrightarrow} x_{1}^{c_1}x_{3}^{wd_3-c_3}(x_{3}^{c_3}-x_{1}^{c_{1}}) \stackrel{g_2} {\longrightarrow} 0.$

For every $0 \leq i \leq w-3$ we have $S(f,q_{i})=x_{1}^{c_1}x_{2}^{(i+2)d_2}-x_{3}^{c_3-(i+2)d_3}x_{4}^{(i+2)c_4} \stackrel{q_{i+1}} {\longrightarrow} 0$. We now consider the $S$--polynomial $S(f,q_{w-2})$. If $c_{3}-wd_{3}>0$, then \[ S(f,q_{w-2})=x_1^{c_1}x_2^{wd_2}-x_3^{c_3-wd_3}x_4^{wc_4}
\stackrel{r}{\longrightarrow}
x_1^{c_1}x_2^{wd_2-c_2}(x_2^{c_2}-x_3^{c_3})
\stackrel{g_1}{\longrightarrow} \] \[
x_1^{c_1}x_2^{wd_2-c_2}(x_1^{c_1}-x_3^{c_3})
\stackrel{g_2}{\longrightarrow}
0.
\]
Suppose that $c_3-wd_3 \leq 0$. Then $S(f,q_{w-2})=x_1^{c_1}x_2^{wd_2}x_3^{wd_3-c_3}-x_4^{wc_4}$. If $c_2-wd_2>0$, then $S(f,q_{w-2})\stackrel{r}{\longrightarrow}0$. If $c_2-wd_2 \leq 0$, then
\[
S(f,q_{w-2})
\stackrel{r}{\longrightarrow}
x_1^{c_1}x_2^{wd_2-c_2}x_3^{wd_3-c_3}(x_2^{c_2}-x_3^{c_3})
\stackrel{g_1}{\longrightarrow} \] \[
x_1^{c_1}x_2^{wd_2-c_2}x_3^{wd_3-c_3}(x_1^{c_1}-x_3^{c_3})
\stackrel{g_2}{\longrightarrow}
0.
\]

Next, we show that $S(p_{i},p_{j}) \stackrel{T}{\longrightarrow} 0$, where $0 \leq i<j \leq w-2$. We compute $S(p_{i},p_{j})=x_{1}^{c_1}x_{3}^{(i+1)d_3}(x_{2}^{(j-i)d_{2}}x_{3}^{(j-i)d_{3}}-x_{4}^{(j-i)c_4}) \stackrel{f}{\longrightarrow} 0$. We now verify that $S(p_i,r) \stackrel{T} {\longrightarrow} 0$ for every $0 \leq i \leq w-2$. Suppose first that $c_2-wd_2>0$. Then $$S(p_i,r)=x_{1}^{c_1}x_{3}^{wd_{3}-c_3}(x_{2}^{c_{2}+(w-i-1)d_{2}}-x_{3}^{c_{3}-(w-i-1)d_{3}}x_{4}^{(w-i-1)c_{4}}) \stackrel{q_{w-i-2}} {\longrightarrow}$$ $$x_{1}^{c_1}x_{2}^{(w-i-1)d_{2}}x_{3}^{wd_{3}-c_3}(x_{2}^{c_2}-x_{1}^{c_{1}}) \stackrel{g_{1}} {\longrightarrow} 0.$$ Now suppose that $c_2-wd_2 \leq 0$. Then $$S(p_i,r)=x_{1}^{c_1}x_{3}^{(i+1)d_3}(x_{2}^{(w-i-1)d_2}x_{3}^{(w-i-1)d_3}-x_{4}^{(w-i-1)c_4}) \stackrel{f} {\longrightarrow} 0.$$

Next, we show that $S(q_{i},q_{j}) \stackrel{T}{\longrightarrow} 0$, where $0 \leq i<j \leq w-2$. We compute $$S(q_{i},q_{j})=x_{1}^{c_1}x_{2}^{(i+1)d_2}(x_{2}^{(j-i)d_{2}}x_{3}^{(j-i)d_{3}}-x_{4}^{(j-i)c_4}) \stackrel{f}{\longrightarrow} 0.$$ We now verify that $S(q_i,r) \stackrel{T} {\longrightarrow} 0$ for every $0 \leq i \leq w-2$. Suppose first that $c_2-wd_2>0$. Then $S(q_i,r)=x_{1}^{c_1}x_{2}^{(i+1)d_2}(x_{2}^{(w-i-1)d_2}x_{3}^{(w-i-1)d_3}-x_{4}^{(w-i-1)c_4}) \stackrel{f} {\longrightarrow}  0$. Now assume that $c_2-wd_2 \leq 0$, so that 
$r=x_4^{wc_4}-x_1^{c_1}x_2^{wd_2-c_2}x_3^{wd_3}$. Then $$S(q_i,r)=x_{1}^{c_1}x_{2}^{wd_{2}-c_2}(x_{3}^{c_3+(w-i-1)d_{3}}-x_{2}^{c_{2}-(w-i-1)d_{2}}x_{4}^{(w-i-1)c_{4}}) \stackrel{p_{w-i-2}} {\longrightarrow}$$ $$x_{1}^{c_1}x_{2}^{wd_{2}-c_2}x_{3}^{(w-i-1)d_{3}}(x_{3}^{c_3}-x_{1}^{c_{1}}) \stackrel{g_{2}} {\longrightarrow} 0.$$

Finally, we show that $S(p_{i},q_{j}) \stackrel{T}{\longrightarrow} 0$, where $0 \leq i, j \leq w-2$. If $i=j$, then 
$S(p_i,q_i)=x_1^{c_1}(x_2^{c_2}-x_3^{c_3}) \stackrel{g_1} {\longrightarrow} x_1^{c_1}(x_1^{c_1}-x_3^{c_3}) \stackrel{g_2} {\longrightarrow} 0$. Suppose that $i<j$. Then 
$$S(p_{i},q_{j})=x_{1}^{c_1}(x_{2}^{c_{2}+(j-i)d_{2}}-x_{3}^{c_{3}-(j-i)d_{3}}x_{4}^{(j-i)c_{4}}) \stackrel{g_{1}}{\longrightarrow}$$ 
$$x_{1}^{c_1}(x_{1}^{c_{1}}x_{2}^{(j-i)d_{2}}-x_{3}^{c_{3}-(j-i)d_{3}}x_{4}^{(j-i)c_{4}}) \stackrel{q_{j-i-1}}{\longrightarrow} 0.$$ Now suppose that $i>j$. Then 
$$S(p_{i},q_{j})=x_{1}^{c_1}(x_{2}^{c_{2}+(j-i)d_{2}}x_{4}^{(i-j)c_{4}}-x_{3}^{c_{3}+(i-j)d_{3}}) \stackrel{p_{i-j-1}}{\longrightarrow}$$ 
$$x_{1}^{c_1}x_{3}^{(i-j)d_{3}}(x_{1}^{c_1}-x_{3}^{c_3}) \stackrel{g_{2}}{\longrightarrow} 0.$$ Thus, $T$ is a Gr\"obner basis for $I({\bf a})$ with respect to $<$. \hfill $\square$ \par\vspace{0.5\baselineskip}


\begin{thm1} \label{CM16} 
Let $d_1=0$, $d_{2}<c_{2}$, $d_{3}<c_{3}$, $c_3>c_2$, and $c_4<d_{2}+d_{3}$. Suppose that the following conditions hold:
\begin{enumerate} 
\item $(w-1)(c_{4}-d_{2}-d_{3})+c_{2}-c_{1} \geq 0$.
\item $w(c_{4}-d_{2}-d_{3})+c_{3}-c_{1} \geq 0$ when $c_{2}-wd_{2}> 0$.
\item $w(c_{4}-d_{2}-d_{3})+c_{2}-c_{1} \geq 0$ when $c_{2}-wd_{2} \leq 0$.
\end{enumerate}
Then $\overline{C({\bf a})}$ is arithmetically Cohen--Macaulay.
\end{thm1}

\noindent \textbf{Proof.} By Proposition \ref{CMBasic12}, the set $T$ is a Gr\"obner basis for $I({\bf a})$ with respect to $<$. Since $x_1$ does not divide the initial monomial of any element in $T$, it follows from \cite[Theorem 2.2]{HerSta} that $\overline{C({\bf a})}$ is arithmetically Cohen--Macaulay. \hfill $\square$


\begin{prop1} \label{ACM7} 
Let $n \geq 3$ be an odd integer, and let $a_1=2n(2n+2)$, $a_2=2n+2$, $a_3=2n$, and $a_4=3n+2$. Then $\overline{C({\bf a})}$ is arithmetically Cohen--Macaulay, and the minimal number of generators of $I^{h}({\bf a})$ is equal to $2n+2$.
\end{prop1}

{\em \noindent Proof.} Since ${\rm gcd}(a_{1}, a_2, a_3)=2$ and $a_4$ is relatively prime to $2$, we deduce that ${\rm gcd}(a_{i} \mid 1 \leq i \leq 4)=1$. Also $I({\bf a})$ is generated by $x_{2}^{2n}-x_1$, $x_{3}^{2n+2}-x_{1}$, and $x_{4}^{2}-x_{2}^{2}x_{3}$. Here, $d_2=2<2n=c_2$, $d_3=1<2n+2=c_3$, $c_4=2$, $w=n$ and $c_2-wd_{2}=0$. Observe that $(w-1)(c_4-d_2-d_3)+c_2-c_1=n>0$ and $w(c_4-d_2-d_3)+c_2-c_1=n-1>0$. Therefore, by Theorem \ref{CM16}, $\overline{C({\bf a})}$ is arithmetically Cohen--Macaulay. Moreover, by Proposition \ref{CMBasic12}, the set $$T=\{x_{2}^{2n}-x_1, x_{3}^{2n+2}-x_{1}, x_{2}^{2}x_3-x_{4}^2\} \cup \{p_{i}=x_{2}^{2n-2(i+1)}x_{4}^{2(i+1)}-x_{1}x_3^{i+1} \mid 0 \leq i \leq n-2\}$$ $$\cup \{q_{i}=x_{3}^{2n+2-(i+1)}x_{4}^{2(i+1)}-x_{1}x_{2}^{2(i+1)} \mid 0 \leq i \leq n-2\} \cup \{r=x_{4}^{2n}-x_{1}x_{3}^{n}\}$$ is a Gr\"obner basis for $I({\bf a})$ with respect to $<$. By Remark \ref{VeryBasicRemark}, the ideal $I^{h}({\bf a})$ is minimally generated by the set $$\{x_{2}^{2n}-x_{0}^{2n-1}x_1, x_{3}^{2n+2}-x_{0}^{2n+1}x_1, x_{2}^{2}x_3-x_{0}x_{4}^{2}\} \cup \{p^{h}_{i} \mid 0 \leq i \leq n-2\} \cup$$ $$\{q^{h}_{i} \mid 0 \leq i \leq n-2\} \cup \{r^{h}\}.$$ Thus, the ideal $I^{h}({\bf a})$ is minimally generated by $2n+2$ elements. \hfill $\square$ \par\vspace{0.5\baselineskip}

{\em Finally, we consider the case $a_{4}={\rm max}\{a_1,\ldots,a_4\}$. Without loss of generality, we fix the assumptions $c_{3}>c_{1}$, $c_{2}>c_{1}$, $d_{2} \leq c_2$, and $d_{3} \leq c_3$.} Under these assumptions, we give a characterization for when $I^{h}({\bf a})$ is a complete intersection.

\begin{thm1} \label{Comp2} 
The ideal $I^{h}({\bf a})$ is a complete intersection if and only if $d_{2}=c_2$ and $d_{3}=c_{3}$.
\end{thm1}

{\em \noindent Proof.}  ($\Leftarrow$) Since $a_{4}={\rm max}\{a_1,\ldots,a_4\}$, we use the degree reverse lexicographic order $<$ with $x_{1}>x_2>x_3>x_4$. Let $f=x_{1}^{d_1+2c_1}-x_{4}^{c_4} \in I({\bf a})$. Then ${\rm in}_{<}(g_1)=x_{2}^{c_2}$, ${\rm in}_{<}(g_2)=x_{3}^{c_3}$, and ${\rm in}_{<}(f)=x_{1}^{d_{1}+2c_1}$ are pairwise relatively prime. Hence, $\{g_1, g_2, f\}$ is a Gr\"obner basis for $I({\bf a})$ with respect to $<$. By Remark \ref{VeryBasicRemark}, the set $\{g_{1}^{h}, g_{2}^{h}, f^{h}\}$ is a Gr\"obner basis for $I^{h}({\bf a})$ with respect to $<_{0}$. Consequently, $I^{h}({\bf a})$ is a complete intersection.\\
($\Rightarrow$) Let $T=\{g_{1}^{h}, g_{2}^{h}, x_{1}^{u_1}-x_{0}^{u_{0}}x_{2}^{u_2}x_{3}^{u_3}x_{4}^{u_4}\}$ be a minimal generating set of $I^{h}({\bf a})$, where $u_{1} \geq c_{1}$ and $u_{1} \geq u_2+u_3+u_4$. Since $a_{4}={\rm max}\{a_1,\ldots,a_4\}$, we have $c_{4}<d_{1}+d_{2}+d_{3}$, and therefore $x_{0}^{d_{1}+d_{2}+d_{3}-c_{4}}x_{4}^{c_{4}}-x_{1}^{d_{1}}x_{2}^{d_{2}}x_{3}^{d_{3}} \in I^{h}({\bf a})$. Since the set $T$ generates $I^{h}({\bf a})$, the monomial $x_{0}^{u_{0}}x_{2}^{u_2}x_{3}^{u_3}x_{4}^{u_4}$ divides $x_{0}^{d_{1}+d_{2}+d_{3}-c_{4}}x_{4}^{c_{4}}$. Thus, $u_{2}=u_{3}=0$ and $u_{4} \leq c_{4}$. Hence, $x_{1}^{u_1}-x_{0}^{u_{0}}x_{4}^{u_4} \in I^{h}({\bf a})$, so $x_{1}^{u_1}-x_{4}^{u_4} \in I({\bf a})$, and therefore $u_{4} \geq c_{4}$. This implies $u_{4}=c_{4}$. Since at least one of $d_2$ and $d_3$ is nonzero, we may assume without loss of generality that $d_{2} \neq 0$. Observe that $x_{1}^{c_1+d_1}x_{3}^{d_3}-x_{2}^{c_{2}-d_2}x_{1}^{u_1} \in I({\bf a})$. As $c_2-d_2<c_2$, we conclude that $u_{1} \geq c_1+d_{1}$, and therefore $x_{3}^{d_3}-x_{1}^{u_1-c_1-d_1}x_{2}^{c_2-d_2} \in I({\bf a})$. Thus $d_{3} \geq c_3$, and since by assumption $d_3 \leq c_3$, we conclude that $d_{3}=c_3$. Now consider the binomial $x_{1}^{c_1+d_1}x_{2}^{d_2}-x_{3}^{c_3-d_3}x_{4}^{c_4} \in I({\bf a})$. Since $c_3=d_3$ and $x_{1}^{u_1}-x_{4}^{c_4} \in I({\bf a})$, we have $x_{1}^{c_1+d_1}x_{2}^{d_2}-x_{1}^{u_1} \in I({\bf a})$, which implies $x_{2}^{d_2}-x_{1}^{u_1-c_1-d_1} \in I({\bf a})$. Therefore, $d_2 \geq c_2$. Since $d_2 \leq c_2$ by assumption, we conclude that $d_2=c_2$. \hfill $\square$ \par\vspace{0.5\baselineskip}

The next example demonstrates the complete intersection property of $I^h(\mathbf{a})$ when $d_2=c_2$ and $d_3=c_3$.

\begin{ex1} 
{\rm Let ${\bf a}=(140, 84, 60, 245)$. Then $I({\bf a})$ is generated by the binomials $x_{2}^{5}-x_{1}^{3}$, $x_{3}^{7}-x_{1}^3$, and $x_{4}^4-x_1x_2^5x_3^7$. Here, $c_{2}=d_{2}=5$ and $c_{3}=d_{3}=7$. By Theorem \ref{Comp2}, the ideal $I^{h}({\bf a})$ is a complete intersection.}
\end{ex1}


\section{Projective closures in case~3}

In this section, we study Case 3 of Theorem \ref{CasesMinimal}, where $I({\bf a})$ has a minimal generating set $\{x_{2}^{c_2}-x_{1}^{c_1}, x_{3}^{c_3}-x_{1}^{b_1}x_{2}^{b_2}, x_{4}^{c_4}-x_{1}^{d_1}x_{2}^{d_2}x_{3}^{d_3}\}$. Both $b_1$ and $b_2$ are nonzero, and at least two of $d_1$, $d_2$, and $d_3$ are nonzero.

Suppose first that $a_{i}={\rm max}\{a_1,\ldots,a_4\}$ for some $i \in \{1,2\}$. Without loss of generality, we assume $a_{1}={\rm max}\{a_1,\ldots,a_4\}$ and $b_2 \leq c_2$. The next theorem provides a complete intersection criterion for $I^{h}({\bf a})$ when $b_2<c_2$.

\begin{thm1} \label{Basic5} 
Suppose that $b_2<c_2$. Then $I^{h}({\bf a})$ is a complete intersection if and only if $c_{3} \geq b_{1}+b_{2}$ and the element $c_{4}a_{4} \in \mathbb{N}\{a_1,\ldots,a_4\}$ admits a factorization $c_{4}a_{4}=v_{1}a_{1}+v_{2}a_{2}+v_{3}a_{3}$ with $c_{4} \geq v_1+v_{2}+v_{3}$ for some $v_1, v_2, v_3 \in \mathbb{N}$. 
\end{thm1}

{\em \noindent Proof.} ($\Leftarrow$) Let $g_{1}=x_{2}^{c_2}-x_{1}^{c_1}, g_{2}=x_{3}^{c_3}-x_{1}^{b_1}x_{2}^{b_2}$, and $g_{3}=x_{4}^{c_4}-x_{1}^{v_1}x_{2}^{v_2}x_{3}^{v_3} \in I({\bf a})$, where $c_{4} \geq v_1+v_{2}+v_{3}$. Consider the degree reverse lexicographic order $<$ with $x_{4}>x_{3}>x_{2}>x_{1}$. Then ${\rm in}_{<}(g_1)=x_{2}^{c_2}$, ${\rm in}_{<}(g_2)=x_{3}^{c_3}$, and ${\rm in}_{<}(g_{3})=x_{4}^{c_4}$ are pairwise relatively prime. Hence, $\{g_{i} \mid 1 \leq i \leq 3\}$ is a Gr\"obner basis for $I({\bf a})$ with respect to $<$. By Remark \ref{VeryBasicRemark}, $\{g_{i}^{h} \mid 1 \leq i \leq 3\}$ is a Gr\"obner basis for $I^{h}({\bf a})$ with respect to $<_{0}$. Therefore, $I^{h}({\bf a})$ is a complete intersection.\\
($\Rightarrow$) Assume first that $c_3<b_1+b_2$. Since $\overline{C({\bf a})}$ is arithmetically Cohen--Macaulay and ${\rm in}_{<}(g_2)=x_1^{b_1}x_2^{b_2}$, \cite[Theorem~2.2]{HerSta} yields a binomial $B\in I({\bf a})$ whose initial monomial divides $x_2^{b_2}$, contradicting the minimality of $c_2$. Now suppose that for every factorization $c_{4}a_{4}=v_{1}a_{1}+v_{2}a_{2}+v_{3}a_{3}$ we have $c_{4}<v_1+v_{2}+v_{3}$. Let $T=\{g_{1}^{h}, g_{2}^{h}, x_{4}^{z_4}-x_{0}^{z_{0}}x_{1}^{z_1}x_{2}^{z_2}x_{3}^{z_3}\}$ be a minimal generating set of $I^{h}({\bf a})$, where $z_{4} \geq z_{1}+z_{2}+z_{3}$. Since $x_{4}^{z_4}-x_{1}^{z_1}x_{2}^{z_2}x_{3}^{z_3} \in I({\bf a})$, we have $z_{4} \geq c_4$, and $c_4<z_1+z_2+z_3$ implies $z_4>c_4$. Moreover, $x_{0}^{d_1+d_2+d_3-c_4}x_{4}^{c_4}-x_{1}^{d_1}x_2^{d_2}x_{3}^{d_3} \in I^{h}({\bf a})$, and since $T$ generates $I^{h}({\bf a})$, we have $x_{4}^{z_4}$ divides $x_{4}^{c_4}$, hence $z_{4} \leq c_4$, a contradiction. \hfill $\square$ \par\vspace{0.5\baselineskip}

The failure of either condition in Theorem \ref{Basic5} prevents $I^{h}({\bf a})$ from being a complete intersection. The following example illustrates this situation.

\begin{ex1} 
{\rm Let ${\bf a}=(a_1=12, a_2=4, a_3=10, a_4=11)$. Then $I({\bf a})$ is generated by $x_{2}^{3}-x_1$, $x_{3}^{2}-x_{1}x_{2}^{2}$, and $x_{4}^{2}-x_{2}^{3}x_{3}$. Here, $b_{2}=2<3=c_{2}$, $c_{3}=2<3=b_{1}+b_{2}$, and $2a_4=a_1+a_3$. By Theorem \ref{Basic5}, $I^{h}({\bf a})$ is not a complete intersection.}
\end{ex1}

Next, we provide a complete intersection criterion for $I^{h}({\bf a})$ when $b_2=c_2$.

\begin{thm1} \label{Basic23} 
Suppose that $b_2=c_2$. Then $I^{h}({\bf a})$ is a complete intersection if and only if the element $c_{4}a_{4} \in \mathbb{N}\{a_1,\ldots,a_4\}$ admits a factorization $c_{4}a_{4}=v_{1}a_{1}+v_{2}a_{2}+v_{3}a_{3}$ for some $v_1, v_2, v_3 \in \mathbb{N}$ with $c_{4} \geq v_1+v_{2}+v_{3}$. 
\end{thm1}

{\em \noindent Proof.} ($\Leftarrow$) Let $g_{1}=x_{2}^{c_2}-x_{1}^{c_1}, g_{2}=x_{3}^{c_3}-x_{1}^{b_1}x_{2}^{b_2}$, and $g_{3}=x_{4}^{c_4}-x_{1}^{v_1}x_{2}^{v_2}x_{3}^{v_3} \in I({\bf a})$, where $c_{4} \geq v_1+v_{2}+v_{3}$. Consider the degree reverse lexicographic order $<$ with $x_{4}>x_{3}>x_{2}>x_{1}$. Then ${\rm in}_{<}(g_1)=x_{2}^{c_2}$, ${\rm in}_{<}(g_2)=x_{3}^{c_3}$, and ${\rm in}_{<}(g_3)=x_{4}^{c_4}$ are pairwise relatively prime. Hence, $\{g_{i} \mid 1 \leq i \leq 3\}$ is a Gr\"obner basis for $I({\bf a})$ with respect to $<$. By Remark \ref{VeryBasicRemark}, $\{g_{i}^{h} \mid 1 \leq i \leq 3\}$ is a Gr\"obner basis for $I^{h}({\bf a})$ with respect to $<_{0}$. Thus, $I^{h}({\bf a})$ is a complete intersection.\\ 
($\Rightarrow$) The proof of this direction follows the same argument as the $(\Rightarrow)$ direction of Theorem~\ref{Basic5}, using only the obstruction arising from the factorization of $c_4a_4$, and is therefore omitted. \hfill $\square$ \par\vspace{0.5\baselineskip}

The following example illustrates Theorem~\ref{Basic23}.

\begin{ex1} 
{\rm Let ${\bf a}=(a_{1}=18, a_{2}=9, a_{3}=12, a_{4}=14)$. Then $I({\bf a})$ is generated by $x_{2}^{2}-x_{1}$, $x_{3}^{3}-x_{1}x_{2}^{2}$, and $x_{4}^{3}-x_{2}^2x_{3}^{2}$. Here, $b_{2}=c_{2}=2$ and $3a_4=a_1+2a_3$. By Theorem \ref{Basic23}, $I^{h}({\bf a})$ is a complete intersection.} 
\end{ex1}

Suppose now that $a_{3}={\rm max}\{a_1,\ldots,a_4\}$. Without loss of generality, we fix the assumption $c_{2}>c_{1}$. The next theorem establishes a complete intersection criterion in the case where $c_4 \geq d_1+d_2+d_3$.

\begin{thm1} \label{basicvfff} 
Suppose that $c_{4} \geq d_1+d_2+d_3$. Then $I^{h}({\bf a})$ is a complete intersection if and only if the element $c_{3}a_3 \in \mathbb{N}\{a_1,\ldots,a_4\}$ admits a factorization $c_{3}a_3=u_1a_{1}$ for some $u_1 \in \mathbb{N}$ with $u_{1}>c_1$.
\end{thm1}

{\em \noindent Proof.} ($\Leftarrow$) Let $f_{1}=x_{2}^{c_2}-x_{1}^{c_1}, f_{2}=x_{4}^{c_4}-x_{1}^{d_1}x_2^{d_2}x_{3}^{d_3}$, and $f_{3}=x_{1}^{u_1}-x_{3}^{c_3} \in I({\bf a})$. Consider the degree reverse lexicographic order $<$ with $x_4>x_2>x_1>x_3$. Since $c_{4} \geq d_1+d_2+d_3$, we have ${\rm in}_{<}(f_{2})=x_{4}^{c_4}$. Then ${\rm in}_{<}(f_{1})=x_{2}^{c_2}$, ${\rm in}_{<}(f_{2})$, and ${\rm in}_{<}(f_{3})=x_{1}^{u_1}$ are pairwise relatively prime. Thus $\{f_i \mid 1 \leq i \leq 3\}$ is a Gr\"obner basis for $I({\bf a})$ with respect to $<$. By Remark \ref{VeryBasicRemark}, $\{f_{i}^h \mid 1 \leq i \leq 3\}$ is a Gr\"obner basis for $I^{h}({\bf a})$ with respect to $<_{0}$. Thus, $I^{h}({\bf a})$ is a complete intersection.

($\Rightarrow$) Let $T=\{x_{1}^{u_1}-x_{0}^{u_0}x_{2}^{u_2}x_{3}^{u_3}x_{4}^{u_4}, x_{2}^{c_2}-x_{0}^{c_2-c_1}x_{1}^{c_1}, x_{4}^{c_4}-x_{0}^{c_4-d_1-d_2-d_3}x_{1}^{d_1}x_{2}^{d_2}x_{3}^{d_3}\}$ be a minimal generating set of $I^{h}({\bf a})$, where $u_{1} \geq u_2+u_3+u_4$. Since $a_{3}={\rm max}\{a_1,\ldots,a_4\}$, we have $c_{3}<b_{1}+b_{2}$, and therefore $x_{0}^{b_{1}+b_{2}-c_{3}}x_{3}^{c_3}-x_{1}^{b_1}x_{2}^{b_2} \in I^{h}({\bf a})$. Furthermore, since $T$ generates $I^{h}({\bf a})$, at least one of $x_{0}^{u_0}x_{2}^{u_2}x_{3}^{u_3}x_{4}^{u_4}$ and $x_{0}^{c_4-d_1-d_2-d_3}x_{1}^{d_1}x_{2}^{d_2}x_{3}^{d_3}$ divides $x_{0}^{b_{1}+b_{2}-c_{3}}x_{3}^{c_3}$. We distinguish two cases.\\ 
(1) Suppose that $x_{0}^{u_0}x_{2}^{u_2}x_{3}^{u_3}x_{4}^{u_4}$ divides $x_{0}^{b_{1}+b_{2}-c_{3}}x_{3}^{c_3}$. Then $u_{2}=u_{4}=0$ and $u_{3} \leq c_{3}$. Hence, $x_{1}^{u_1}-x_{0}^{u_0}x_{3}^{u_3} \in I^{h}({\bf a})$, implying $x_{1}^{u_1}-x_{3}^{u_3} \in I({\bf a})$. Thus $u_{3} \geq c_{3}$, hence $u_{3}=c_{3}$. Since $c_{1}a_1 \neq c_{3}a_3$, we conclude that $u_1>c_1$. Furthermore, the binomial $x_{0}^{d_1+d_2+d_3-c_4}x_{4}^{c_4}-x_{1}^{d_1}x_{2}^{d_2}x_{3}^{d_3}$ belongs to $I^{h}({\bf a})$. Since $T$ generates $I^{h}({\bf a})$, the monomial $x_{4}^{v_4}$ must divide $x_{0}^{d_{1}+d_{2}+d_3-c_{4}}x_{4}^{c_4}$, hence $v_{4} \leq c_{4}$. But $x_{4}^{v_{4}}-x_{1}^{v_{1}}x_{2}^{v_{2}}x_{3}^{v_{3}} \in I({\bf a})$, so $v_{4} \geq c_{4}$, and thus $v_{4}=c_{4}$. Consequently, condition (A) holds.\\
(2) Suppose that $x_{0}^{v_{0}}x_{1}^{v_{1}}x_{2}^{v_{2}}x_{4}^{v_{4}}$ divides $x_{0}^{b_{1}+b_{2}-c_{3}}x_{3}^{c_3}$. Then $v_{1}=v_{2}=0$ and $v_{3} \leq c_{3}$, so $x_{4}^{v_4}-x_{0}^{v_0}x_{3}^{v_3} \in I^{h}({\bf a})$, implying $x_{4}^{v_4}-x_{3}^{v_3} \in I({\bf a})$. Thus $v_{3} \geq c_{3}$, hence $v_{3}=c_{3}$. Since $c_{3}a_{3} \neq c_4a_4$, it follows that $v_{4}>c_{4}$. Next, consider the binomial $x_{0}^{b_{1}+b_{2}-c_{3}}x_{3}^{c_{3}}-x_{1}^{b_{1}}x_{2}^{b_{2}}$. Since $T$ generates $I^{h}({\bf a})$, the monomial $x_{0}^{u_0}x_{2}^{u_2}x_{3}^{u_3}x_{4}^{u_4}$ must divide $x_{0}^{b_{1}+b_{2}-c_{3}}x_{3}^{c_{3}}$, so $u_2=u_4=0$ and $u_3 \leq c_{3}$. Hence, $x_{1}^{u_1}-x_{0}^{u_0}x_{3}^{u_3} \in I^{h}({\bf a})$ with $u_1 \geq u_3$, implying $x_{1}^{u_1}-x_{3}^{u_3} \in I({\bf a})$ with $u_3 \geq c_{3}$. Thus $u_3=c_{3}$. Since $a_1 \neq a_3$, it follows that $u_1>c_{3}$, and because $c_1a_1 \neq c_{3}a_3$, we also have $u_1>c_{1}$. Thus condition (B) holds.  \hfill $\square$

\begin{rem1} {\rm We do not know whether conditions {\rm (A)} and {\rm (B)} in Theorem~\ref{Basic26} can hold simultaneously.}
\end{rem1}

We conclude this section with an example illustrating condition {\rm (A)} in Theorem~\ref{Basic26}.

\begin{ex1} 
{\rm Let ${\bf a}=(a_{1}=3553,a_2=2431,a_3=3069,a_{4}=5168)$. Then $I({\bf a})$ is generated by $x_{2}^{19}-x_{1}^{13}$, $x_{3}^{17}-x_{1}x_{2}^{20}$, and $x_{4}^{11}-x_{1}^{3}x_{2}^{19}$. Here, $c_3=17<21=b_1+b_2$ and $c_4=11$. Since $c_4a_4$ does not admit a factorization $c_4a_4=v_3a_3$ with $v_3>c_3$, condition (B) of Theorem \ref{Basic26} fails. However, since $16a_1=11a_4$ and $17a_3=14a_1+a_2$, condition (A) of Theorem \ref{Basic26} is satisfied. Thus, $I^{h}({\bf a})$ is a complete intersection.}
\end{ex1}


